% \iffalse meta-comment
%
% Copyright (C) 1993-2026
% The LaTeX Project and any individual authors listed elsewhere
% in this file.
%
% This file is part of the LaTeX base system.
% -------------------------------------------
%
% It may be distributed and/or modified under the
% conditions of the LaTeX Project Public License, either version 1.3c
% of this license or (at your option) any later version.
% The latest version of this license is in
%    https://www.latex-project.org/lppl.txt
% and version 1.3c or later is part of all distributions of LaTeX
% version 2008 or later.
%
% This file has the LPPL maintenance status "maintained".
%
% The list of all files belonging to the LaTeX base distribution is
% given in the file `manifest.txt'. See also `legal.txt' for additional
% information.
%
% The list of derived (unpacked) files belonging to the distribution
% and covered by LPPL is defined by the unpacking scripts (with
% extension .ins) which are part of the distribution.
%
% \fi
% \iffalse
%%% From File: ltfloat.dtx
%
%<*driver>
% \fi
\ProvidesFile{ltfloat.dtx}[2026-10-06 v1.2k LaTeX Kernel (Floats)]
% \iffalse
\documentclass{ltxdoc}
\GetFileInfo{ltfloat.dtx}
\title{\filename}
\date{\filedate}
 \author{%
  Johannes Braams\and
  David Carlisle\and
  Alan Jeffrey\and
  Leslie Lamport\and
  Frank Mittelbach\and
  Chris Rowley\and
  Rainer Sch\"opf}

\begin{document}
 \MaintainedByLaTeXTeam{latex}
 \maketitle
 \DocInput{\filename}
\end{document}
%</driver>
% \fi
%
% \providecommand\hook[1]{\texttt{#1}}
% \providecommand\plug[1]{\texttt{#1}}
%
% \section{Floats}
%
%  The different types of floats are identified by a \meta{type} name,
%  which is the name of the counter for that kind of float.  For
%  example, figures are of type `figure' and tables are of type `table'.
%  Each \meta{type} has associated a positive \meta{type number}, which
%  is a power of two.  E.g.,\\
%  figures might be have type number~1, tables type number~2, programs
%  type number~4, etc.
%
%  The locations where a float can go are specified by a
%  \meta{placement specifier}, which is a list of the possible
%  locations, each denoted by a letter as follows:
%    \begin{center}
%    \begin{tabular}{l@{ : }l@{ --- }l}
%     h & here   & at the current location in the text.\\
%     t & top    & at the top of a text page.\\
%     b & bottom & at the bottom of a text page.\\
%     p & page   & on a separate float page
%    \end{tabular}
%    \end{center}
%  In addition, in conjunction with these, you can use `!' which means
%  that the current values of the float positioning parameters are
%  ignored for this float. (Has no effect on `p', float page
%  positioning.)
%  For example, `pht' specifies that the float can appear in any of
%  three locations: page, here or top.
%
% \MaybeStop{}
%
%
% \changes{v1.0a}{1994/03/04}{Initial version, split from latex.dtx}
% \changes{v1.0b}{1994/03/28}{Split further from ltherest.dtx}
% \changes{v1.0e}{1994/04/25}{Changed warning messages}
% \changes{v1.0e}{1994/04/25}{Removed obsolete tracing code}
% \changes{v1.0f}{1994/05/03}
%     {(CAR) Added \cs{@largefloatcheck}}
% \changes{v1.0f}{1994/05/03}{Removed unnecessary braces from
% arguments of \cs{@ifnextchar}}
% \changes{v1.0i}{1994/05/22}{Use new warning commands}
% \changes{v1.1e}{1994/11/17}
%         {\cs{@tempa} to \cs{reserved@a}}
% \changes{v1.1g}{1994/12/10}{Some temps reinserted temporarily}
% \changes{v1.1n}{1995/11/28}{documentation fixes}
% \changes{v1.1s}{1997/06/16}{documentation fixes}
%
% \subsection{Floating Environments}
%    \begin{macrocode}
%<*2ekernel>
\message{floats,}
%    \end{macrocode}
% \begin{oldcomments}
%
% Where floats may appear on a page, and how many may appear there
% are specified by the following float placement parameters.  The
% numbers are named like counters so the user can set them with
% the ordinary counter-setting commands.
%
%  \c@topnumber      : Number of floats allowed at the top of a column.
%  \topfraction      : Fraction of column that can be devoted to floats.
%  \c@dbltopnumber, \dbltopfraction
%                    : Same as above, but for double-column floats.
%  \c@bottomnumber, \bottomfraction
%                    : Same as above for bottom of page.
%  \c@totalnumber    : Number of floats allowed in a single column,
%                          including in-text floats.
%  \textfraction     :Minimum fraction of column that must contain text.
%  \floatpagefraction: Minimum fraction of page that must be taken
%                          up by float page.
%  \dblfloatpagefraction
%                    : Same as above, for double-column floats.
%
% The document style must define the following.
%
%    \fps@TYPE   : The default placement specifier for floats of type
%                  TYPE.
%
%    \ftype@TYPE : The type number for floats of type TYPE.
%
%    \ext@TYPE   : The file extension indicating the file on which the
%                  contents list for float type TYPE is stored.
%                    For example,  \ext@figure = 'lof'.
%
%    \fnum@TYPE  : A macro to generate the figure number for a caption.
%                  For example, \fnum@TYPE == Figure \thefigure.
%
%    \@makecaption{NUM}{TEXT} :
%              A macro to make a caption, with NUM the value
%              produced by \fnum@... and TEXT the text of the caption.
%              It can assume it's in a \parbox of the appropriate width.
%
% \@float{TYPE}[PLACEMENT] : This macro begins a float environment for a
%     single-column float of type TYPE with PLACEMENT as the placement
%     specifier.  The default value of PLACEMENT is defined by
%     \fps@TYPE.  The environment is ended by \end@float.
%     E.g., \figure == \@float{figure}, \endfigure == \end@float.
%
%  \@float{TYPE}[PLACEMENT] ==
%   BEGIN
%     if hmode then \@bsphack
%                   \@floatpenalty := -10002
%              else \@floatpenalty := -10003
%     fi
%     \@captype ==L TYPE
%     \@dblflset
%     \@fps     ==L PLACEMENT
%     \@onelevel@sanitize \@fps
%     add default PLACEMENT if at most ! in PLACEMENT == \@fpsadddefault
%     if inner
%       then LaTeX Error: 'Not in outer paragraph mode.'
%            \@floatpenalty := 0
%       else if \@freelist nonempty
%              then \@currbox  :=L head of \@freelist
%                   \@freelist :=G tail of \@freelist
%                   \count\@currbox :=G 32*\ftype@TYPE +
%                                          bits determined by PLACEMENT
%              else \@floatpenalty := 0
%                   LaTeX Error: 'Too many unprocessed floats'
%            fi
%     fi
%     \@currbox :=G   \color@vbox
%                       \normalcolor
%                         \vbox{
%                          %% 15 Dec 87 --
%                          %% removed \boxmaxdepth :=L 0pt
%                          %% that made box 0 depth because it screwed
%                          %% things up. Instead, added \vskip0pt at end
%                               \hsize = \columnwidth
%                               \@parboxrestore
%                               \@floatboxreset
%   END
%
%  \caption ==
%    BEGIN
%     \refstepcounter{\@captype}
%     \@dblarg{\@caption{\@captype}}
%    END
%
% In following definition, \par moved from after \addcontentsline to
% before \addcontentsline because the \write could cause
% an extra blank line to be added to the paragraph above the
% caption.  (Change made 12 Jun 87)
%
%  \@caption{TYPE}[STEXT]{TEXT} ==
%   BEGIN
%     \par
%     \addcontentsline{\ext@TYPE}{TYPE}{\numberline{\theTYPE}{STEXT}}
%     \begingroup
%       \@parboxrestore
%       \@normalsize
%       \@makecaption{\fnum@TYPE}{TEXT}
%       \par
%     \endgroup
%   END
%
%
%  \@dblfloat{TYPE}[PLACEMENT] : Macro to begin a float environment for
%     a double-column float of type TYPE with PLACEMENT as the placement
%     specifier.  The default value of PLACEMENT is 'tp'
%     The environment is ended by \end@dblfloat.
%     E.g., \figure* == \@dblfloat{figure},
%           \endfigure* == \end@dblfloat.
%
%  \@dblfloat{TYPE}[PLACEMENT] ==
%     Identical to \@float{TYPE}[PLACEMENT] except \hsize and \linewidth
%     are set to \textwidth.
% \end{oldcomments}
%
% \begin{macro}{\@floatpenalty}
%    \begin{macrocode}
\newcount\@floatpenalty
%    \end{macrocode}
% \end{macro}
%
% \paragraph{The caption interface.}
%
% \changes{v1.2k}{2026-10-06}{Caption interface: hook \texttt{caption/before},
%    socket \texttt{caption/step}, switches \cs{if@captionstep} and
%    \cs{if@captiontarget}}
% \changes{v1.2k}{2026-10-06}{Hook \texttt{caption/prepare}, socket
%    \texttt{caption/listentry}}
% \changes{v1.2k}{2026-10-06}{Socket \texttt{caption/typeset}}
%    \cs{caption}\oarg{list text}\marg{text} and
%    \cs{caption*}\oarg{opt}\marg{text} process a caption in the
%    following steps (inside a float, \meta{type} is the expansion of
%    \cs{@captype}, e.g., \texttt{figure}):
%    \begin{enumerate}
%    \item The generic hook \hook{cmd/caption/before} (as for every
%      command; it is executed before anything else).
%    \item Outside of a float an error is raised (as before).
%    \item \cs{@kernel@caption@reset} undoes what an earlier caption in
%      the same group may have left (see below) and sets the switches
%      \cs{if@captionstar}, \cs{if@captionstep} and \cs{if@captiontarget}
%      to their defaults. Then the star and the arguments are read. An
%      absent optional argument is \cs{NoValue}.
%    \item The hook \hook{caption/before} is executed.
%    \item Unless \cs{caption*} was used: the socket \texttt{caption/step}
%      steps the counter.
%    \item \cs{@caption}\marg{type}\oarg{list text}\marg{text} is called
%      (the list text is the optional argument, or the text if it was
%      not given). It does a \cs{par}, executes the hook
%      \hook{caption/prepare}, writes the list entry with the socket
%      \texttt{caption/listentry} (not for \cs{caption*}), typesets the
%      caption with the socket \texttt{caption/typeset} and resets the
%      switches \cs{if@captionstar}, \cs{if@captionstep} and
%      \cs{if@captiontarget} to their defaults.
%    \end{enumerate}
%    Code that calls \cs{@caption} directly (e.g., for sub-captions) gets
%    what \cs{@caption} does; it can set \cs{@captionstartrue} before the
%    call.
%    Without additional code in the hooks and with the default plugs, the
%    result is the same as in earlier releases,
%    except that \cs{caption*} is available.
%
% \paragraph{\cs{caption*} with classes and packages.}
%    What \cs{caption*} gives depends on the definitions that are active:
%    \begin{itemize}
%    \item With the kernel \cs{caption} and \cs{@caption} (also if a
%      package wraps them and calls the kernel definitions, as
%      \cls{memoir} and \pkg{nameref} do): no step, no target, no list
%      entry, and \cs{@makecaption} is called with the label
%      \cs{@caption@nolabel} while \cs{if@captionstar} is true.
%    \item With the kernel \cs{caption} and an \cs{@caption} that a class
%      or package has replaced by a definition that does not know the
%      star (\cls{llncs}, \cls{mwart}, \ldots): the kernel
%      lets \cs{fnum@}\meta{type} be \cs{@caption@nolabel}, makes
%      \cs{ext@}\meta{type} empty (so \cs{addcontentsline} writes to no
%      list) and gives \cs{theH}\meta{type} a unique suffix (so a target
%      that the definition makes has a new name), until the next
%      \cs{caption} or the end of the group (see
%      \cs{@caption@nostep@defs}). Such a definition, if it uses these
%      commands as the kernel did, then also gives a caption without
%      label, without list entry and without a duplicate target, in the
%      layout of the class, also if it defers the typesetting
%      (\pkg{threeparttable} does that for a caption above the table).
%      \cs{if@captionstar} is true until the same point. For a caption
%      that does not step its counter for another reason
%      (\cs{@captionstepfalse}, a continued float) only the unique
%      suffix is used.
%    \item The class's \cs{@makecaption} decides how the caption looks.
%      The contract is: \cs{if@captionstar} is true, and the label (the
%      first argument) is \cs{@caption@nolabel}, never empty. A
%      \cs{@makecaption} that supports \cs{caption*} tests the switch and
%      omits the label \emph{and} its separator (the tagging code does
%      that); this is the only documented test (a test for an empty label
%      does not work). For the many definitions that do not test it,
%      \cs{@caption@nolabel} typesets nothing and removes the usual
%      separators after it (a colon or a period and the following
%      space, also across the end of a group or of \cs{textbf}), so the
%      standard classes, many classes derived from them and classes such
%      as \cls{llncs}, \cls{IEEEtran} and \cls{mwart} show neither label
%      nor separator without any change. A separator in a macro
%      (\pkg{KOMA-Script}, \cls{memoir}, \cls{amsart}, \pkg{babel-french}
%      for French with pdf\TeX), another space (\cs{hskip},
%      \cs{hspace}) or another character is printed before the text; in
%      most cases a warning is given then (once), see
%      \cs{@caption@nolabel}. The standard classes are not changed, because packages
%      compare \cs{@makecaption} with their definition (for example
%      \pkg{babel-french}, which replaces the separator for French only if
%      \cs{@makecaption} is the standard one).
%    \item A package or class that replaces \cs{caption} itself bypasses
%      all of this, for example \pkg{hyperref} if there is no
%      \cs{DocumentMetadata}, \cls{beamer} and some classes:
%      their \cs{caption} takes the star as the text of a numbered
%      caption, as before. Such code has to be changed (for example by
%      keeping the kernel \cs{caption} and using the hooks and sockets
%      described here); the kernel cannot do it, because it is not called.
%      For \pkg{float}, which is no longer maintained, the first aid lets
%      \cs{caption*} use the kernel definition.
%    \end{itemize}
%
%    The declarations are also written to \texttt{latexrelease}, guarded,
%    so that rolling forward from an older format works. After a roll-back
%    they stay defined but are not used.
%
% \paragraph{The hooks \texttt{caption/before} and \texttt{caption/prepare}.}
%    Two hooks with four arguments each:
%    \begin{description}
%    \item[\texttt{caption/before}] is executed by \cs{caption} once per
%      caption (also for \cs{caption*}), after the star and the arguments
%      have been read and before the counter is stepped. The arguments are
%      the type, \cs{BooleanTrue} for \cs{caption*} or \cs{BooleanFalse}
%      otherwise (test with \cs{IfBooleanTF}), the optional argument as
%      given or \cs{NoValue} if it was absent (test with \cs{IfNoValueTF};
%      |[]| and |[{}]| are both given and empty), and the caption text.
%      The hook is executed at the group level of the caller (in a float
%      the group of the float, with its \cs{linewidth}), so local settings
%      made here are valid for the rest of the float. This is the place to
%      decide that the counter should not be stepped (\cs{@captionstepfalse},
%      e.g., for a continued float) or that no target should be created
%      (\cs{@captiontargetfalse}, if the code makes its own target). The hook
%      is not executed if \cs{@caption} is called directly.
%    \item[\texttt{caption/prepare}] is executed by \cs{@caption} after the
%      \cs{par} and before the list entry, for every call of \cs{@caption}
%      (also for \cs{caption*} and for direct calls). The arguments are the
%      type, the star state as above (from \cs{if@captionstar}), the list
%      text and the caption text. At this point \cs{@currentlabel} is set
%      and, if the step made a target, \cs{@currentHref} as well;
%      \cs{@currentlabelname} is the list text (not for \cs{caption*}). This is
%      the place for code that today wraps \cs{@caption}, e.g., to make a
%      target of its own or to record the title of the caption for named
%      references.
%    \end{description}
%    Both hooks are executed outside of any group the kernel makes.
%    \begin{macrocode}
%</2ekernel>
%<*2ekernel|latexrelease>
%<latexrelease>\IncludeInRelease{2026/11/01}%
%<latexrelease>                 {\@kernel@caption@step}{caption interface}%
\NewHookWithArguments{caption/before}{4}
\NewHookWithArguments{caption/prepare}{4}
%    \end{macrocode}
%
% \begin{socketdecl}{caption/step}
%    This socket is used by \cs{caption} (not by \cs{caption*}) to step
%    the counter of the caption and to set \cs{@currentlabel} and the
%    target for links (\cs{@currentHref}). The argument is the counter
%    (the type of the caption). It is used after the hook
%    \hook{caption/before}. Other code may use it for caption-like
%    steps (e.g., for sub-captions or a phantom caption); it then has to
%    set the switches \cs{if@captionstep} and \cs{if@captiontarget}
%    itself before the use.
%
%    Every plug must fulfil the following contract:
%    \begin{enumerate}
%    \item Everything is wrapped in the socket \texttt{refstepcounter}
%      (as in \cs{refstepcounter}).
%    \item If \cs{if@captionstep} is true, the counter is stepped with
%      \cs{stepcounter} (so that the counters reset by it are reset).
%      Otherwise no counter is changed.
%    \item \cs{@currentcounter} and \cs{@currentlabel} are set as by
%      \cs{refstepcounter}.
%    \item If \cs{if@captiontarget} is true: if \cs{@floatHref@}\meta{type}
%      is defined, not \cs{relax} and not empty, \cs{@currentHref} is
%      set globally to its value, and no new target is made. Otherwise
%      the sockets \texttt{refstepcounter/target} and \texttt{recordtarget}
%      are used, as in \cs{refstepcounter}; if the counter was not stepped,
%      \cs{theH}\meta{counter} gets a unique suffix |*|\meta{n} for the
%      socket \texttt{refstepcounter/target} (locally). If
%      \cs{if@captiontarget} is false,
%      none of this is done and \cs{@currentHref} is unchanged.
%    \end{enumerate}
%    A target made from the unchanged counter would have the same name as
%    the one of the caption that stepped the counter last, so a caption
%    that does not step the counter (e.g., the caption of a continued
%    float) must never get such a target: \pkg{hyperref} would make a
%    second destination with the same name, and links to it would lead to
%    the first caption. With the suffix the name is unique (for example
%    |figure.1*1|; \pkg{hyperref} does not use |*| in
%    \cs{theH}\meta{counter}), so a \cs{label} after such a caption gives
%    the number of the counter and a link to the caption itself.
%    \meta{n} counts the targets made in this way for the same name, so
%    the targets of the continued parts of figure~1 are |figure.1*1|,
%    |figure.1*2|, \ldots, and adding a continued float elsewhere does not
%    change these names. The suffix is added to the definition of
%    \cs{theH}\meta{counter} without expanding it first, so \pkg{hyperref}
%    expands it in its usual context; if \cs{theH}\meta{counter} is not
%    defined, \cs{the}\meta{counter} is used (as \pkg{hyperref} does). The
%    target is made by the same socket as for other captions, so
%    \pkg{hyperref} makes it as usual and without \pkg{hyperref} no
%    node is added to the list. (\cs{NextLinkTarget} is not used for the
%    name: the kernel plug of \texttt{refstepcounter/target} does not
%    honour it, so without \pkg{hyperref} \cs{@currentHref} and the
%    target recorded for tagging would still be the name of the last
%    stepped caption.) Inside a float that has a target of
%    its own (\cs{@floatHref@}\meta{type}), that target is used, because it
%    belongs to the float that contains the caption.
%    The default plug \plug{kernel} uses \cs{@kernel@caption@step}, which
%    implements exactly this.
% \end{socketdecl}
% \begin{plugdecl}{kernel (caption/step)}
%    \begin{macrocode}
\IfSocketExistsF{caption/step}{%
  \NewSocket{caption/step}{1}%
  \NewSocketPlug{caption/step}{kernel}{\@kernel@caption@step{#1}}%
  \AssignSocketPlug{caption/step}{kernel}%
}
%    \end{macrocode}
% \end{plugdecl}
%
% \begin{socketdecl}{caption/listentry}
%    This socket writes the entry for the list of figures, tables, etc.
%    It is used by \cs{@caption} after the hook \hook{caption/prepare},
%    unless the star form \cs{caption*} was used (\cs{if@captionstar}).
%    The first argument is the caption type, the second the text for the
%    list (the optional argument of \cs{caption}, or the caption text if
%    it was not given), the third the caption text. The default plug
%    \plug{kernel} writes the entry with \cs{addcontentsline} as before;
%    it also writes an (empty) entry if the optional argument is empty.
%    A package that wants to know whether the optional argument was
%    given can record that in the hook \hook{caption/before}.
% \end{socketdecl}
% \begin{plugdecl}{kernel (caption/listentry)}
%    \begin{macrocode}
\IfSocketExistsF{caption/listentry}{%
  \NewSocket{caption/listentry}{3}%
  \NewSocketPlug{caption/listentry}{kernel}
    {%
      \addcontentsline{\csname ext@#1\endcsname}{#1}%
        {\protect\numberline{\csname the#1\endcsname}{\ignorespaces #2}}%
    }%
  \AssignSocketPlug{caption/listentry}{kernel}%
}
%    \end{macrocode}
% \end{plugdecl}
%
% \begin{socketdecl}{caption/typeset}
%    This socket typesets the caption. It is used by \cs{@caption} for
%    every caption, after the list entry. The arguments are the caption
%    type, \cs{BooleanTrue} for \cs{caption*} or \cs{BooleanFalse}
%    otherwise, and the caption text (without \cs{ignorespaces}).
%
%    A plug is used at the group level of the caller of \cs{@caption}
%    (in a float the group of the float, with its \cs{linewidth}), outside
%    of any group the kernel makes, and \cs{@parboxrestore} has not been
%    executed yet. It has to do its own grouping; settings that it makes
%    outside of its group are valid after the caption, so it should only
%    make them on purpose. \cs{if@captionstar} is valid while it runs.
%    A plug that typesets the caption with \cs{@makecaption} should pass
%    \cs{@caption@nolabel} as first argument for \cs{caption*}, so that
%    classes and the tagging code behave as with the default plug.
%
%    The default plug \plug{kernel} uses \cs{@kernel@caption@typeset}. A
%    package that typesets captions completely (e.g., \pkg{caption}) or
%    that has to redirect them can assign its own plug; there is one
%    owner, the last assignment wins. A plug can fall back to
%    \cs{@kernel@caption@typeset}.
% \end{socketdecl}
% \begin{plugdecl}{kernel (caption/typeset)}
%    \begin{macrocode}
\IfSocketExistsF{caption/typeset}{%
  \NewSocket{caption/typeset}{3}%
  \NewSocketPlug{caption/typeset}{kernel}
    {\@kernel@caption@typeset{#1}{#2}{#3}}%
  \AssignSocketPlug{caption/typeset}{kernel}%
}
%    \end{macrocode}
% \end{plugdecl}
%
% \begin{macro}{\if@captionstar,\if@captionstep,\if@captiontarget}
%    \cs{if@captionstar} is true while a caption made by \cs{caption*} is
%    processed: from the start of \cs{caption*} until the end of
%    \cs{@caption}, so also in the hooks, in the plugs of
%    \texttt{caption/typeset} and in \cs{@makecaption}. In that case the first
%    argument of \cs{@makecaption} is \cs{@caption@nolabel} (or, if
%    \cs{@caption} has been replaced, \cs{fnum@}\meta{type}, which is then
%    \cs{@caption@nolabel}), and a \cs{@makecaption} can test this switch
%    to omit the label and its separator. If \cs{@caption} is not
%    the kernel's, the switch stays true until the next \cs{caption} or
%    the end of the current group, as the settings of
%    \cs{@caption@nostep@defs}, because that \cs{@caption} may defer its
%    work. Such an \cs{@caption} that tests the switch itself (or whose
%    \cs{@makecaption} does) should therefore set it false at its end;
%    otherwise a later \cs{@makecaption} in the same group (for example
%    for the caption of a listing) would omit its label as well. Code
%    that calls \cs{@caption} directly can set it true before the
%    call.
%
%    \cs{if@captionstep} and \cs{if@captiontarget} are true outside of
%    captions. They are set true by \cs{caption} before the hook
%    \hook{caption/before} is executed. Code in that hook (or code that
%    uses the socket \texttt{caption/step} itself, directly before the use)
%    can set them false:
%    \begin{itemize}
%    \item \cs{@captionstepfalse}: the counter is not stepped by the plug,
%      either because the caption continues an earlier one (a continued
%      float) or because the code has already stepped it itself (e.g.,
%      for a caption after sub-captions that stepped it);
%      \cs{@currentcounter} and \cs{@currentlabel}
%      are still set, and if \cs{if@captiontarget} is true a target as
%      well: the float target or one with a unique name, never a second
%      target with the name made from the counter (see the socket
%      \texttt{caption/step}). The plug cannot tell the two cases apart,
%      so in both the target gets the unique name; code that has stepped
%      the counter itself and wants the usual name sets
%      \cs{@captiontargetfalse} as well and makes the target itself
%      (with the socket \texttt{refstepcounter/target}).
%    \item \cs{@captiontargetfalse}: the step makes no target and does not
%      record one for tagging, and \cs{@currentHref} is unchanged; the code
%      makes its own target.
%    \end{itemize}
%    The three switches are reset at the end of \cs{@caption} (with
%    \cs{@kernel@caption@reset}), so nothing is passed on to a later
%    caption.
%
%    (\cs{@ifundefined} and not \cs{ifx} because a skipped \cs{if...}
%    would unbalance the conditional.)
%    \begin{macrocode}
\@ifundefined{if@captionstar}{\newif\if@captionstar}{}
\@ifundefined{if@captionstep}{\newif\if@captionstep\@captionsteptrue}{}
\@ifundefined{if@captiontarget}{\newif\if@captiontarget\@captiontargettrue}{}
%    \end{macrocode}
% \end{macro}
%
% \begin{macro}{\@kernel@caption@step,\@kernel@caption@unique}
% \begin{macro}{\@caption@target,\@caption@newtarget,\@caption@floattarget,
%                \@caption@uniquetarget,\@caption@uniquedo,\@caption@uniqueH,
%                \@caption@uniqueH@from,\@caption@Htype,\@caption@noH}
%    \cs{@kernel@caption@step}\marg{counter} implements the contract of
%    the socket \texttt{caption/step}. It is not \cs{refstepcounter},
%    because the counter is not always stepped and the target not always
%    made. The star guard of \cs{refstepcounter} is not needed, caption
%    counters are never |*|.
%    \begin{macrocode}
\def\@kernel@caption@step#1{%
  \UseSocket{refstepcounter}{%
    \if@captionstep
      \stepcounter{#1}%
    \fi
    \edef\@currentcounter{#1}%
    \protected@edef\@currentlabel
       {\csname p@#1\expandafter\endcsname\csname the#1\endcsname}%
    \if@captiontarget
      \@caption@target{#1}%
    \fi}}
%    \end{macrocode}
%    The float code can provide the target of the current float of a
%    type in \cs{@floatHref@}\meta{type} (the code for tagged floats in
%    \pkg{latex-lab-float} does that). If it is defined, not \cs{relax}
%    and not empty, the caption uses this target, otherwise a new one is
%    made as by \cs{refstepcounter}; if the counter was not stepped,
%    \cs{@caption@uniquetarget} makes it with a unique name (with
%    \cs{@kernel@caption@unique}), because a target named after the
%    unchanged counter would be a duplicate. \cs{@ifundefined}
%    is used and not
%    \cs{ifcsname}, so that a name that has been made \cs{relax} by a
%    \cs{csname} test counts as undefined. As the type is part of the name,
%    a \texttt{table} caption inside a \texttt{figure} or a step of a
%    sub-caption counter (e.g., \texttt{subfigure}) does not use the target
%    of the figure.
% \changes{v1.2k}{2026-10-06}{A caption that does not step the counter
%    gets a target with a unique name (suffix |*|\meta{n}, counted per
%    name)}
%    \begin{macrocode}
\def\@caption@target#1{%
  \@ifundefined{@floatHref@#1}%
    {\@caption@newtarget{#1}}%
    {\expandafter\ifx\csname @floatHref@#1\endcsname\@empty
       \expandafter\@caption@newtarget
     \else
       \expandafter\@caption@floattarget
     \fi
     {#1}}}
\def\@caption@newtarget#1{%
  \if@captionstep
    \expandafter\@firstoftwo
  \else
    \expandafter\@secondoftwo
  \fi
  {\UseSocket{refstepcounter/target}{#1}}%
  {\@caption@uniquetarget{#1}}%
  \UseTaggingSocket{recordtarget}}
\def\@caption@uniquetarget#1{%
  \@caption@uniquedo{#1}{\UseSocket{refstepcounter/target}{#1}}}
%    \end{macrocode}
% \changes{v1.2k}{2026-10-06}{\cs{@kernel@caption@unique} added}
%    \cs{@kernel@caption@unique}\marg{counter}\marg{code} executes
%    \meta{code} while \cs{theH}\meta{counter} gives the unique name of a
%    caption that does not step its counter: |*|\meta{n} is appended to
%    it, with the next \meta{n} for this name. The numbers are the same
%    as those of the kernel's own targets for such captions, so a name
%    made in this way and a kernel target never collide. It is
%    documented for code that makes the target of a caption itself or
%    only needs its name (for example the \pkg{caption} package, which
%    sets \cs{@captiontargetfalse} and makes its anchor in the hook
%    \hook{caption/prepare} with
%    |\@kernel@caption@unique{|\meta{type}|}{\hyper@makecurrent{|\meta{type}|}}|),
%    so that it does not have to know how the names are counted:
%    \begin{itemize}
%    \item Every use takes a new number. Code that wants one name for
%      one caption uses it once and keeps the result (e.g.,
%      \cs{@currentHref}).
%    \item In the hooks \hook{caption/before} and \hook{caption/prepare}
%      and in the plugs of the caption sockets (and outside of captions)
%      \cs{theH}\meta{counter} has its normal meaning, and no number is
%      reserved, so the numbers have no gaps. (A number that the kernel
%      reserves for a replaced \cs{@caption}, see
%      \cs{@caption@nostep@defs}, is given back at the start of the kernel
%      \cs{@caption}, before \hook{caption/prepare}.)
%    \item After a \cs{caption*} or a caption that does not step, if
%      \cs{@caption} is a replaced one, \cs{theH}\meta{counter} has the
%      unique form of that caption until the next \cs{caption} or the end
%      of the group (see \cs{@caption@nostep@defs}). The command can also
%      be used there: the name is made from the normal meaning, which is
%      saved, so it is the next unique name after that of the caption.
%    \item The command adds no group: assignments in \meta{code} survive,
%      also local ones (such as \cs{@currentHref} with the \pkg{hyperref}
%      option \texttt{localanchorname}). Afterwards
%      \cs{theH}\meta{counter} has its previous meaning again (or is
%      undefined again). \meta{code} must not use the command for the
%      same counter.
%    \end{itemize}
%    The code is \cs{@caption@uniquedo}, which the kernel uses itself;
%    \cs{@kernel@caption@unique} is defined as it together with
%    \cs{@kernel@caption} (see there), so that after a roll-back with
%    \pkg{latexrelease} to an earlier date it is undefined, as the
%    interface it belongs to.
%    \begin{macrocode}
\long\def\@caption@uniquedo#1#2{%
  \ifcsname theH#1\endcsname
    \expandafter\let\csname @caption@unique@#1\expandafter\endcsname
      \csname theH#1\endcsname
  \else
    \expandafter\let\csname @caption@unique@#1\endcsname\@caption@noH
  \fi
  \@caption@uniqueH{#1}%
  #2%
  \expandafter\ifx\csname @caption@unique@#1\endcsname\@caption@noH
    \expandafter\let\csname theH#1\endcsname\@undefined
  \else
    \expandafter\let\csname theH#1\expandafter\endcsname
      \csname @caption@unique@#1\endcsname
  \fi}
\def\@caption@noH{\@caption@noH}
%    \end{macrocode}
%    \cs{@caption@uniqueH}\marg{counter} (internal; use
%    \cs{@kernel@caption@unique}) redefines \cs{theH}\meta{counter}
%    locally as its previous meaning (kept in
%    \cs{@caption@H@}\meta{counter}, or \cs{the}\meta{counter} if
%    \cs{theH}\meta{counter} is undefined) followed by |*|\meta{n}. The
%    previous meaning is not expanded here, so \pkg{hyperref} expands it
%    in the context it sets up for names. \meta{n} is counted per name:
%    \cs{@caption@Hname} is the name without the suffix (made with
%    \cs{protected@edef}, only to tell names apart), and
%    \cs{@caption@H@n@}\meta{name} holds the last \meta{n} used for it,
%    globally, so the numbering does not depend on captions elsewhere.
%
%    While the settings of \cs{@caption@nostep@defs} for the same counter
%    are active (\cs{@caption@nostep@restore} is not \cs{relax}),
%    \cs{theH}\meta{counter} is already the unique form and refers to
%    \cs{@caption@H@}\meta{counter}; taking it as the previous meaning
%    would let that command refer to itself (an endless loop when the name
%    is expanded). Then the meaning saved by \cs{@caption@nostep@defs}
%    (\cs{@caption@saved@H}) is used. For \cs{@caption@nostep@defs} itself
%    it is the same as the current one.
%    \begin{macrocode}
\def\@caption@uniqueH#1{%
  \edef\@caption@Htype{#1}%
  \ifx\@caption@nostep@restore\relax
    \@caption@uniqueH@from{theH#1}{#1}%
  \else
    \ifx\@caption@Htype\@caption@nostep@type
      \@caption@uniqueH@from{@caption@saved@H}{#1}%
    \else
      \@caption@uniqueH@from{theH#1}{#1}%
    \fi
  \fi
  \protected@edef\@caption@Hname{#1.\csname @caption@H@#1\endcsname}%
  \edef\@caption@Hname{\detokenize\expandafter{\@caption@Hname}}%
  \@ifundefined{@caption@H@n@\@caption@Hname}%
    {\expandafter\xdef\csname @caption@H@n@\@caption@Hname\endcsname{1}}%
    {\expandafter\xdef\csname @caption@H@n@\@caption@Hname\endcsname
       {\the\numexpr\csname @caption@H@n@\@caption@Hname\endcsname+\@ne}}%
  \expandafter\edef\csname theH#1\endcsname
    {\expandafter\noexpand\csname @caption@H@#1\endcsname
     *\csname @caption@H@n@\@caption@Hname\endcsname}}
\def\@caption@uniqueH@from#1#2{%
  \@ifundefined{#1}%
    {\expandafter\let\csname @caption@H@#2\expandafter\endcsname
       \csname the#2\endcsname}%
    {\expandafter\let\csname @caption@H@#2\expandafter\endcsname
       \csname #1\endcsname}}
\def\@caption@floattarget#1{%
  \xdef\@currentHref{\csname @floatHref@#1\endcsname}}
%    \end{macrocode}
% \end{macro}
% \end{macro}
%
% \begin{macro}{\@kernel@caption@typeset}
%    \cs{@kernel@caption@typeset}\marg{type}\marg{star}\marg{text} is the
%    typesetting of the kernel: the previous code of \cs{@caption},
%    with the label \cs{@caption@nolabel} for \cs{caption*}. It is
%    documented, so that a plug of \texttt{caption/typeset} can fall back
%    to it.
%    \begin{macrocode}
\long\def\@kernel@caption@typeset#1#2#3{%
  \begingroup
%    \end{macrocode}
%
% The paragraph setting parameters are normalized at this point, however
% |\@parboxrestore| resets |\everypar| which is not correct in this
% context so |\@setminipage| is called if needed.
%
% The float mechanism, like minipage, sets the flag |@minipage| true
% before executing the user-supplied text. Many \LaTeX\ constructs
% test for this flag and do not add vertical space when it is true.
% The intention is that this emulates \TeX's `top of page' behaviour.
% The flag must be set false at the start of the first paragraph. This
% is achieved by a redefinition of |\everypar|, but the call to
% |\@parboxrestore| removes that redefinition, so it is re-inserted
% if needed. If the flag is already false then the |\caption| was not
% the first entry in the float, and so some other paragraph has already
% activated the special |\everypar|. In this case no further action is
% needed.
%    \begin{macrocode}
    \@parboxrestore
    \if@minipage
      \@setminipage
    \fi
%    \end{macrocode}
%
%    \begin{macrocode}
    \normalsize
    \IfBooleanTF{#2}%
      {\@makecaption{\@caption@nolabel}}%
      {\@makecaption{\csname fnum@#1\endcsname}}%
      {\ignorespaces #3}\par
  \endgroup}
%    \end{macrocode}
% \end{macro}
%
% \begin{macro}{\@caption@nolabel}
% \begin{macro}{\@caption@nolabel@look,\@caption@nolabel@do,
%                \@caption@nolabel@next,\@caption@nolabel@drop,
%                \@caption@nolabel@punct,\@caption@nolabel@space,
%                \@caption@nolabel@group,\@caption@nolabel@stop,
%                \@caption@nolabel@clean,\@caption@nolabel@other,
%                \@caption@nolabel@brace,\@caption@nolabel@line,
%                \@caption@nolabel@hmode,\@caption@nolabel@m,
%                \@caption@nolabel@nowarn,\@caption@nolabel@warn}
% \changes{v1.2k}{2026-10-06}{Macro added: \cs{@caption@nolabel}, the label
%    of \cs{caption*}}
% \changes{v1.2k}{2026-10-06}{The label of \cs{caption*} removes the usual
%    separators after it}
%    \cs{@caption@nolabel} is the label that \cs{caption*} passes to
%    \cs{@makecaption}. The contract for \cs{@makecaption} is only that
%    \cs{if@captionstar} is true and that the label is not empty but
%    typesets nothing; a \cs{@makecaption} that supports \cs{caption*}
%    tests the switch and omits the label and its separator.
%
%    Everything else is a convenience for the many \cs{@makecaption}
%    definitions that do not (yet) test the switch. It is separate from
%    the contract: if \cs{@caption@nolabel@look} were \cs{@gobble}, the
%    label would simply typeset nothing. While \cs{if@captionstar} is true,
%    the label looks at the tokens after it:
%    \begin{itemize}
%    \item A colon or a period directly after the label is removed
%      (spaces before it are skipped), then spaces, |~|,
%      \cs{nobreakspace}, \cs{space}, |\ |, \cs{enspace}, \cs{enskip},
%      \cs{quad} and \cs{qquad}.
%    \item If the label, or what was removed after it, is the last thing
%      in a group or in the argument of \cs{textbf} or another command
%      defined with \cs{DeclareTextFontCommand} (recognized by
%      \cs{check@icr}), this continues after the end of the group (with
%      \cs{aftergroup}), where a colon or a period can still follow if
%      nothing has been removed yet (|{#1}: #2|, |\textbf{#1}: #2|). The
%      italic correction \cs{maybe@ic} that \cs{check@icr} puts after
%      the group is removed, since there is nothing to correct.
%    \item If something was removed, the paragraph is started
%      (\cs{leavevmode}), in the group of the removed material, as that
%      material would have done. This is not done if \cs{par} follows in
%      that group, so a label line such as |{\centering #1.\par}| gives
%      no empty line; a \cs{par} after the end of the group
%      (|{\bfseries #1.}\par|) comes too late and still gives one.
%    \item If, in the group of the label, the label or the removed
%      material is followed by anything else than the text (which starts
%      with \cs{ignorespaces}), \cs{par}, \cs{relax}, \cs{@},
%      \cs{vadjust}, \cs{penalty}, the end of a conditional, the end of a box that
%      measures the label (\cs{sbox},
%      \cs{settowidth}) or, after a removed separator, a group (which can
%      contain the text), it is taken as material that is printed before
%      the text: a separator in a macro (such as \cs{captionformat} of
%      \pkg{KOMA-Script}, \cs{@contdelim} of \cls{memoir} or the
%      separator of \pkg{babel-french}), other space (\cs{hskip}, also in
%      braces) or a character (|--|). Then a warning is given once per
%      document. The paragraph is not started for it (unless something
%      was removed before it), so that in vertical mode a space at the
%      start of a separator macro is still ignored, as without the label,
%      and vertical material (\cs{vskip}) stays vertical. After the end
%      of a group, the class code can continue with anything (for
%      example with code that measures or boxes the label), so nothing is
%      checked there: a separator after a group that contains only the
%      label (|\textit{#1}\textbf{:}|, |\fbox{#1}:|) is printed without a
%      warning. \cs{relax} also ends the look-ahead without a warning, so
%      a separator after it is printed silently as well
%      (|#1\relax\hskip1em|, and |#1\protect\sep| while typesetting, where
%      \cs{protect} is \cs{relax}).
%    \item \cs{unskip}, \cs{unkern} and \cs{unpenalty} directly after the
%      label or the removed material are removed as well: they are meant
%      for the end of the label, which is empty, and in vertical mode
%      they would remove vertical material instead (such as
%      \cs{abovecaptionskip}).
%    \item The rule for the paragraph: it is started before the
%      primitives whose meaning depends on the mode (\cs{kern},
%      \cs{penalty}, \cs{hbox}, \cs{vbox}, \cs{vtop}, \cs{box},
%      \cs{copy}, \cs{raise}, \cs{lower}, \cs{vadjust}) and before
%      \verb|\\|, \cs{newline} and \cs{@}, so that they act in horizontal
%      mode, as after a label, and do not give an error in vertical mode.
%      It is not started before \cs{par}, \cs{vskip}, a space or a macro
%      (macros are not looked into). \cs{kern}, the boxes, \cs{raise},
%      \cs{lower}, \verb|\\| and \cs{newline} also warn, as they print
%      space, material or an empty line before the text; \cs{penalty},
%      \cs{vadjust} and \cs{@} do not.
%    \end{itemize}
%    Nothing of the text is removed, because the text passed to
%    \cs{@makecaption} starts with \cs{ignorespaces}. All of this is
%    only done while \cs{if@captionstar} is true, so the label does
%    nothing else if it ends up somewhere else.
%    The command is robust, so that it survives \cs{edef} and case
%    changing of the label.
%    \begin{macrocode}
\protected\def\@caption@nolabel{%
  \if@captionstar
    \expandafter\@caption@nolabel@look\expandafter P%
  \fi}
%    \end{macrocode}
%    The argument of \cs{@caption@nolabel@look} is the state, a letter:
%    |P| (nothing removed yet, in the group of the label), |Q| (nothing
%    removed, after the end of a group), |S| (something removed, in its
%    group) or |G| (something removed, after the end of its group). A
%    colon or a period is only removed in states |P| and |Q|, and only
%    |P| and |S| can warn. \cs{kernel@ifnextchar} skips (removes)
%    spaces and leaves the next token in \cs{@let@token}. Then
%    \cs{@caption@nolabel@do} decides what to do with that token. To
%    keep conditionals balanced, \cs{@let@token}, which can be \cs{fi},
%    is only used directly after \cs{ifx} and never in a branch that
%    may be skipped.
%    \begin{macrocode}
\def\@caption@nolabel@look#1{%
  \kernel@ifnextchar\relax
    {\@caption@nolabel@do#1}{\@caption@nolabel@do#1}}
\def\@caption@nolabel@do#1{%
  \def\@caption@nolabel@next{\@caption@nolabel@stop#1}%
  \ifx\@let@token\egroup\@caption@nolabel@group#1\fi
  \ifx\@let@token\check@icr\@caption@nolabel@group#1\fi
  \ifx\@let@token\maybe@ic\@caption@nolabel@drop#1\fi
  \ifx\@let@token\unskip\@caption@nolabel@drop#1\fi
  \ifx\@let@token\unkern\@caption@nolabel@drop#1\fi
  \ifx\@let@token\unpenalty\@caption@nolabel@drop#1\fi
  \ifx\@let@token:\@caption@nolabel@punct#1\fi
  \ifx\@let@token.\@caption@nolabel@punct#1\fi
  \ifx\@let@token~\@caption@nolabel@space#1\fi
  \ifx\@let@token\nobreakspace\@caption@nolabel@space#1\fi
  \ifx\@let@token\space\@caption@nolabel@space#1\fi
  \ifx\@let@token\ \@caption@nolabel@space#1\fi
  \ifx\@let@token\enspace\@caption@nolabel@space#1\fi
  \ifx\@let@token\enskip\@caption@nolabel@space#1\fi
  \ifx\@let@token\quad\@caption@nolabel@space#1\fi
  \ifx\@let@token\qquad\@caption@nolabel@space#1\fi
  \@caption@nolabel@next}
\def\@caption@nolabel@drop#1{%
  \def\@caption@nolabel@next##1{\@caption@nolabel@look#1}}
\def\@caption@nolabel@punct#1{%
  \if P#1\@caption@nolabel@drop S\fi
  \if Q#1\@caption@nolabel@drop S\fi}
\def\@caption@nolabel@space#1{%
  \if G#1\@caption@nolabel@drop G\else\@caption@nolabel@drop S\fi}
\def\@caption@nolabel@group#1{%
  \edef\@caption@nolabel@next{%
    \if S#1\noexpand\leavevmode\fi
    \aftergroup\noexpand\@caption@nolabel@look
    \aftergroup\if P#1Q\else\if S#1G\else#1\fi\fi}}
%    \end{macrocode}
%    At any other token the label is done. \cs{par} and \cs{@@par} end
%    without a paragraph. The tokens that are no separator
%    (\cs{@caption@nolabel@clean}) only start the paragraph if
%    something was removed; a group is such a token only after a
%    removed separator (\cs{@caption@nolabel@brace}). The end of a conditional is recognized by
%    its meaning as a string, so that no \cs{fi} appears in a test.
%    \cs{penalty}, \cs{vadjust} and \cs{@} always start the paragraph
%    (\cs{@caption@nolabel@hmode}); the other mode-dependent primitives,
%    \verb|\\| and \cs{newline} also warn (\cs{@caption@nolabel@line}).
%    Anything else (\cs{@caption@nolabel@other}) is printed: in states
%    |P| and |S| this warns; only in state |S| it starts the paragraph.
%    \begin{macrocode}
\def\@caption@nolabel@stop#1{%
  \let\@caption@nolabel@next\@caption@nolabel@other
  \ifx\@let@token\ignorespaces
    \let\@caption@nolabel@next\@caption@nolabel@clean\fi
  \ifx\@let@token\bgroup\let\@caption@nolabel@next\@caption@nolabel@brace\fi
  \ifx\@let@token\relax\let\@caption@nolabel@next\@caption@nolabel@clean\fi
  \ifx\@let@token\@\let\@caption@nolabel@next\@caption@nolabel@hmode\fi
  \ifx\@let@token\vadjust\let\@caption@nolabel@next\@caption@nolabel@hmode\fi
  \ifx\@let@token\penalty\let\@caption@nolabel@next\@caption@nolabel@hmode\fi
  \ifx\@let@token\kern\let\@caption@nolabel@next\@caption@nolabel@line\fi
  \ifx\@let@token\hbox\let\@caption@nolabel@next\@caption@nolabel@line\fi
  \ifx\@let@token\vbox\let\@caption@nolabel@next\@caption@nolabel@line\fi
  \ifx\@let@token\vtop\let\@caption@nolabel@next\@caption@nolabel@line\fi
  \ifx\@let@token\box\let\@caption@nolabel@next\@caption@nolabel@line\fi
  \ifx\@let@token\copy\let\@caption@nolabel@next\@caption@nolabel@line\fi
  \ifx\@let@token\raise\let\@caption@nolabel@next\@caption@nolabel@line\fi
  \ifx\@let@token\lower\let\@caption@nolabel@next\@caption@nolabel@line\fi
  \ifx\@let@token\color@endgroup
    \let\@caption@nolabel@next\@caption@nolabel@clean\fi
  \ifx\@let@token\ResumeTagging
    \let\@caption@nolabel@next\@caption@nolabel@clean\fi
  \ifx\@let@token\par\let\@caption@nolabel@next\@gobble\fi
  \ifx\@let@token\@@par\let\@caption@nolabel@next\@gobble\fi
  \ifx\@let@token\\\let\@caption@nolabel@next\@caption@nolabel@line\fi
  \ifx\@let@token\newline\let\@caption@nolabel@next\@caption@nolabel@line\fi
  \edef\@caption@nolabel@m{,\meaning\@let@token,}%
  \@expandtwoargs\in@\@caption@nolabel@m\@caption@nolabel@nowarn
  \ifin@\let\@caption@nolabel@next\@caption@nolabel@clean\fi
  \@caption@nolabel@next#1}
\edef\@caption@nolabel@nowarn{,\string\fi,\string\else,\string\or,}
\def\@caption@nolabel@clean#1{\if S#1\expandafter\leavevmode\fi}
\def\@caption@nolabel@other#1{%
  \if\if Q#1G\else#1\fi G\else
    \@caption@nolabel@warn
  \fi
  \@caption@nolabel@clean#1}
\def\@caption@nolabel@hmode#1{\leavevmode}
\def\@caption@nolabel@brace#1{%
  \if P#1\expandafter\@caption@nolabel@other
  \else\expandafter\@caption@nolabel@clean\fi#1}
\def\@caption@nolabel@line#1{\@caption@nolabel@other#1\leavevmode}
\def\@caption@nolabel@warn{%
  \global\let\@caption@nolabel@warn\relax
  \@latex@warning{%
    \string\caption* is used with a caption layout of the\MessageBreak
    class or of a package that does not support it,\MessageBreak
    so a separator or a space may be printed before\MessageBreak
    the text%
  }}
%    \end{macrocode}
% \end{macro}
% \end{macro}
%
%    \begin{macrocode}
%</2ekernel|latexrelease>
%<latexrelease>\EndIncludeInRelease
%<*2ekernel>
%    \end{macrocode}
%
% \begin{macro}{\caption}
% \begin{macro}{\@caption@args,\@caption@read,\@caption@process}
% \begin{macro}{\@caption@nostep@defs,\@caption@nostep@restore,
%                \@caption@nostep@giveback,\@caption@uniqueH@back,
%                \@caption@nostep@type,\@kernel@caption@reset}
% \begin{macro}{\@caption@skip,\@caption@skip@star,\@caption@skip@args,
%                \@caption@skip@opt}
%
%    This is set to be an error message outside a float since no
%    captype is defined there; this may need to be changed by some
%    classes. After the error the star and the arguments are skipped
%    (before, the caption text was typeset). \cs{@caption@skip} removes
%    the rest of \cs{caption} (as \cs{@gobble} did before) and
%    \cs{@caption@skip@star} is protected, so that \cs{caption} outside a
%    float still expands without errors in an expansion context.
% \changes{v1.1u}{1999/04/19}
%    {Made caption an error outside a float: latex/2815}
% \changes{v1.2k}{2026-10-06}
%    {Outside a float, skip the star and the arguments after the error}
% \changes{v1.2k}{2026-10-06}
%    {Add star form; read the arguments without \cs{@dblarg}; hook
%     \texttt{caption/before}; step the counter through the socket
%     \texttt{caption/step}; \cs{@kernel@caption@reset}}
% \changes{v1.2k}{2026-10-06}
%    {Settings for a replaced \cs{@caption} until the next caption}
%
%    The star form \cs{caption*} typesets a caption without a label:
%    it does not step the counter, sets no target and writes no list
%    entry. \cs{@currentlabel} and \cs{@currentHref} are not changed (as
%    after \cs{section*}), so a \cs{label} after \cs{caption*} records
%    whatever was set last, which can be a caption, a section or nothing;
%    it should not be used there. An optional argument is accepted; it
%    is only passed to the
%    hook \hook{caption/before} (and to \cs{@caption} as list text, which
%    the kernel \cs{@caption} only passes to \hook{caption/prepare}).
%
%    The star is read with \cs{@ifstar} and the optional argument with
%    \cs{kernel@ifnextchar}, so spaces before them are skipped, as with
%    \cs{@dblarg} before. The type passed on is the expansion of
%    \cs{@captype}. User text is never inside a conditional that \TeX{}
%    might skip. \cs{@caption} is the last token, as before
%    (\pkg{float}'s \cs{float@caption}, for example, looks for an
%    optional argument after it), also for \cs{caption*}.
%    \begin{macrocode}
%</2ekernel>
%<*2ekernel|latexrelease>
%<latexrelease>\IncludeInRelease{2026/11/01}%
%<latexrelease>                 {\caption}{caption arguments}%
\def\caption{%
   \ifx\@captype\@undefined
     \@latex@error{\noexpand\caption outside float}\@ehd
     \expandafter\@caption@skip
   \else
     \expandafter\@firstofone
   \fi
   {\@kernel@caption@reset
    \@ifstar{\@captionstartrue\@caption@args}\@caption@args}%
}
\long\def\@caption@skip#1{\@caption@skip@star}
\protected\def\@caption@skip@star{%
  \@ifstar\@caption@skip@args\@caption@skip@args}
\def\@caption@skip@args{%
  \kernel@ifnextchar[\@caption@skip@opt{\@caption@skip@opt[]}}
\long\def\@caption@skip@opt[#1]#2{}
\def\@caption@args{%
  \kernel@ifnextchar[\@caption@read{\@caption@read[\NoValue]}}
\long\def\@caption@read[#1]#2{%
  \expandafter\@caption@process\expandafter{\@captype}{#1}{#2}}
\long\def\@caption@process#1#2#3{%
  \if@captionstar
    \expandafter\@firstoftwo
  \else
    \expandafter\@secondoftwo
  \fi
  {\UseHookWithArguments{caption/before}{4}{#1}{\BooleanTrue}{#2}{#3}%
   \@caption@nostep@defs{#1}%
   \IfNoValueTF{#2}%
     {\@caption{#1}[{#3}]{#3}}%
     {\@caption{#1}[{#2}]{#3}}}%
  {\UseHookWithArguments{caption/before}{4}{#1}{\BooleanFalse}{#2}{#3}%
   \UseSocket{caption/step}{#1}%
   \if@captionstep\else
     \@caption@nostep@defs{#1}%
   \fi
   \IfNoValueTF{#2}%
     {\@caption{#1}[{#3}]{#3}}%
     {\@caption{#1}[{#2}]{#3}}}}
%    \end{macrocode}
%    \cs{@caption@nostep@defs}\marg{type} prepares the definitions for a
%    caption that does not step its counter (\cs{caption*}, or
%    \cs{@captionstepfalse} for a continued float), in case \cs{@caption}
%    is not the kernel's. A class or package may have replaced
%    \cs{@caption} by a definition that does not know the star or the
%    switches (\cls{llncs}, \cls{mwart} and others): it would typeset the
%    label with the number of the previous caption, write a list entry
%    and, with \pkg{hyperref}, make a target with the name of the previous
%    caption (\cls{mwart} does that with \cs{hyper@makecurrent}). So:
%    \begin{itemize}
%    \item \cs{theH}\meta{type} gets a unique suffix (see
%      \cs{@caption@uniqueH}), so a target made from it has a new name;
%    \item for \cs{caption*} also \cs{fnum@}\meta{type} is
%      \cs{@caption@nolabel} (the same label as from the kernel, and
%      \cs{if@captionstar} is true) and \cs{ext@}\meta{type} is empty, so
%      that \cs{addcontentsline} writes to no list.
%    \end{itemize}
%    Such a definition, if it uses these commands as the kernel did, then
%    gives a caption without label, without list entry and without a
%    duplicate target, in the layout of the class. The settings are
%    local and last until \cs{@caption@nostep@restore} is executed: at
%    the start of the kernel \cs{@caption} (so its hooks and plugs see
%    the normal definitions), at the start of the next \cs{caption}
%    (with \cs{@kernel@caption@reset}), or at the end of the group; why,
%    is explained below.
%    \cs{addcontentsline} itself is not changed, so an
%    \verb|\addcontentsline{lof}...| after \cs{caption*} works. An
%    undefined \cs{theH}\meta{type} is saved as undefined (\cs{ifcsname}
%    does not make it \cs{relax}), so that it is undefined again after
%    the restore; \cs{@caption@nostep@type} records the type, for
%    \cs{@caption@uniqueH}. The
%    kernel \cs{@caption} also gives the number of the unique suffix
%    back (\cs{@caption@nostep@giveback}), since nothing used it there.
%    \begin{macrocode}
\def\@caption@nostep@defs#1{%
  \edef\@caption@nostep@type{#1}%
  \ifcsname theH#1\endcsname
    \expandafter\let\expandafter\@caption@saved@H\csname theH#1\endcsname
  \else
    \let\@caption@saved@H\@undefined
  \fi
  \if@captionstar
    \expandafter\let\expandafter\@caption@saved@fnum
      \csname fnum@#1\endcsname
    \expandafter\let\expandafter\@caption@saved@ext
      \csname ext@#1\endcsname
    \def\@caption@nostep@restore{%
      \expandafter\let\csname theH#1\endcsname\@caption@saved@H
      \expandafter\let\csname fnum@#1\endcsname\@caption@saved@fnum
      \expandafter\let\csname ext@#1\endcsname\@caption@saved@ext
      \let\@caption@nostep@restore\relax
      \let\@caption@nostep@giveback\relax}%
    \expandafter\let\csname fnum@#1\endcsname\@caption@nolabel
    \expandafter\let\csname ext@#1\endcsname\@empty
  \else
    \def\@caption@nostep@restore{%
      \expandafter\let\csname theH#1\endcsname\@caption@saved@H
      \let\@caption@nostep@restore\relax
      \let\@caption@nostep@giveback\relax}%
  \fi
  \@caption@uniqueH{#1}%
  \edef\@caption@nostep@giveback{%
    \noexpand\@caption@uniqueH@back{\@caption@Hname}%
      {\csname @caption@H@n@\@caption@Hname\endcsname}}}
\let\@caption@nostep@restore\relax
\let\@caption@nostep@giveback\relax
\def\@caption@uniqueH@back#1#2{%
  \expandafter\ifnum\csname @caption@H@n@#1\endcsname=#2\relax
    \expandafter\xdef\csname @caption@H@n@#1\endcsname
      {\the\numexpr#2-\@ne}%
  \fi}
%    \end{macrocode}
%    The settings of \cs{@caption@nostep@defs} have to outlive
%    \cs{caption*} (and a caption that does not step), because the kernel
%    \cs{caption} gets no control back after \cs{@caption}: \cs{@caption}
%    reads its own arguments and has to stay the last token of
%    \cs{caption}, as \pkg{float} and others look for an optional
%    argument after it. And a replaced \cs{@caption} may defer its work:
%    \pkg{threeparttable} stores the call of the previous \cs{@caption}
%    (\cs{TPT@docapt}) and executes it only at the start of the tabular,
%    after \cs{caption*} has ended. Then \cs{fnum@}\meta{type},
%    \cs{ext@}\meta{type}, \cs{theH}\meta{type} and \cs{if@captionstar}
%    must still have the values for \cs{caption*}; the switch also if
%    the deferred \cs{@caption} is the kernel's. So the settings last
%    until the kernel \cs{@caption} starts, until the next \cs{caption},
%    or until the end of the group (in practice that of the float).
%
%    \cs{@kernel@caption@reset} ends them: it undoes everything
%    \cs{@caption@nostep@defs} has set (it drops a pending give-back of
%    the unique number without giving it back, since a replaced
%    \cs{@caption} may have used it for a target), sets
%    \cs{if@captionstar} false and \cs{if@captionstep} and
%    \cs{if@captiontarget} true, which is the state outside of a caption.
%    The kernel \cs{caption} starts with it and the kernel \cs{@caption}
%    ends with it, so it is the only code that ends these settings. It is
%    the command for
%    code that keeps its own \cs{caption} and passes only \cs{caption*}
%    to the kernel (for example with
%    \verb|\@ifstar{\@kernel@caption*}{|\ldots\verb|}|, as the first aid
%    for \pkg{float} does): such code must start its own numbered
%    \cs{caption} with \cs{@kernel@caption@reset}, otherwise the settings
%    of a \cs{caption*} before it in the same float are still active (no
%    label, no list entry, the target name of the \cs{caption*}). As
%    whatever \cs{@caption@nostep@defs} may set in the future is undone
%    by it as well, such code depends on nothing else. It is defined
%    together with \cs{caption} (also by \pkg{latexrelease}), so after a
%    roll-back to an earlier date it is undefined, as \cs{@kernel@caption}.
%    \begin{macrocode}
\def\@kernel@caption@reset{%
  \@caption@nostep@restore
  \@captionstarfalse
  \@captionsteptrue
  \@captiontargettrue}
%</2ekernel|latexrelease>
%<latexrelease>\EndIncludeInRelease
%<latexrelease>\IncludeInRelease{0000/00/00}%
%<latexrelease>                 {\caption}{caption arguments}%
%<latexrelease>\def\caption{%
%<latexrelease>   \ifx\@captype\@undefined
%<latexrelease>     \@latex@error{\noexpand\caption outside float}\@ehd
%<latexrelease>     \expandafter\@gobble
%<latexrelease>   \else
%<latexrelease>     \refstepcounter\@captype
%<latexrelease>     \expandafter\@firstofone
%<latexrelease>   \fi
%<latexrelease>   {\@dblarg{\@caption\@captype}}%
%<latexrelease>}
%<latexrelease>\let\@kernel@caption@reset\@undefined
%<latexrelease>\EndIncludeInRelease
%<*2ekernel>
%    \end{macrocode}
% \end{macro}
% \end{macro}
% \end{macro}
% \end{macro}
%
% \begin{macro}{\@caption}
% \begin{macro}{\@caption@body}
% \changes{v1.0b}{1994/03/28}
%     {Use \cs{normalsize} not \cs{@normalsize}}
% \changes{v1.1r}{1996/12/06}
%     {Call \cs{@setminpage} if needed. latex/2318}
% \changes{v1.2k}{2026-10-06}
%     {No list entry and the label \cs{@caption@nolabel} for \cs{caption*};
%      end with \cs{@kernel@caption@reset}}
% \changes{v1.2k}{2026-10-06}
%     {Undo the settings for a replaced \cs{@caption} before the \cs{par}}
% \changes{v1.2k}{2026-10-06}
%     {Hook \texttt{caption/prepare}; list entry with the socket
%      \texttt{caption/listentry}}
% \changes{v1.2k}{2026-10-06}
%     {Typesetting with the socket \texttt{caption/typeset}}
% \changes{v1.2k}{2026-10-06}
%     {Store the list text in \cs{@currentlabelname}}
%    \cs{@caption}\marg{type}\oarg{list text}\marg{text} keeps its
%    signature, so code that wraps or replaces it keeps working.
%    The typesetting (with \cs{@parboxrestore} and \cs{@makecaption})
%    is now done by the plug of \texttt{caption/typeset}, see
%    \cs{@kernel@caption@typeset}. At the start the settings of
%    \cs{@caption@nostep@defs} are undone, before the \cs{par}, so that
%    the paragraph hooks, the output routine and the hooks and plugs of
%    the caption see the normal definitions. At the end
%    \cs{@kernel@caption@reset} resets the switches (the settings are
%    already undone by then), so that code which calls \cs{@caption} or
%    the socket \texttt{caption/step} directly (without \cs{caption}) gets
%    a normal caption.
%
%    For named references, the list text of a caption (not of
%    \cs{caption*}) is stored in \cs{@currentlabelname}, so that a
%    \cs{label} after the caption records it as the title (as for
%    headings), and not the title of the last heading.
%    \begin{macrocode}
%</2ekernel>
%<*2ekernel|latexrelease>
%<latexrelease>\IncludeInRelease{2026/11/01}%
%<latexrelease>                 {\@caption}{caption body}%
\long\def\@caption#1[#2]#3{%
  \@caption@nostep@giveback
  \@caption@nostep@restore
  \par
  \if@captionstar
    \expandafter\@firstoftwo
  \else
    \expandafter\@secondoftwo
  \fi
  {\@caption@body{#1}{\BooleanTrue}}%
  {\@caption@body{#1}{\BooleanFalse}}%
  {#2}{#3}%
  \@kernel@caption@reset}
\long\def\@caption@body#1#2#3#4{%
  \IfBooleanF{#2}{\def\@currentlabelname{#3}}%
  \UseHookWithArguments{caption/prepare}{4}{#1}{#2}{#3}{#4}%
  \IfBooleanF{#2}{\UseSocket{caption/listentry}{#1}{#3}{#4}}%
  \UseSocket{caption/typeset}{#1}{#2}{#4}}
%</2ekernel|latexrelease>
%<latexrelease>\EndIncludeInRelease
%<latexrelease>\IncludeInRelease{0000/00/00}%
%<latexrelease>                 {\@caption}{caption body}%
%<latexrelease>\long\def\@caption#1[#2]#3{%
%<latexrelease>  \par
%<latexrelease>  \addcontentsline{\csname ext@#1\endcsname}{#1}%
%<latexrelease>    {\protect\numberline{\csname the#1\endcsname}{\ignorespaces #2}}%
%<latexrelease>  \begingroup
%<latexrelease>    \@parboxrestore
%<latexrelease>    \if@minipage
%<latexrelease>      \@setminipage
%<latexrelease>    \fi
%<latexrelease>    \normalsize
%<latexrelease>    \@makecaption{\csname fnum@#1\endcsname}{\ignorespaces #3}\par
%<latexrelease>  \endgroup}
%<latexrelease>\EndIncludeInRelease
%<*2ekernel>
%    \end{macrocode}
% \end{macro}
% \end{macro}
%
% \begin{macro}{\@kernel@caption,\@kernel@caption@unique,\@kernel@@caption}
% \changes{v1.2k}{2026-10-06}{Macro \cs{@kernel@caption} added}
% \changes{v1.2k}{2026-10-06}{Macro \cs{@kernel@caption@unique} added}
% \changes{v1.2k}{2026-10-06}{Macro \cs{@kernel@@caption} added}
%    \cs{@kernel@caption} is a copy of the kernel definition of
%    \cs{caption} (the naming follows \cs{@kernel@refstepcounter}). Code
%    that keeps its own \cs{caption} can pass \cs{caption*} to it (with
%    \verb|\@ifstar{\@kernel@caption*}{|\ldots\verb|}|), and code that
%    has replaced \cs{caption} can restore it.
%    \cs{@kernel@caption@unique} is the documented name of
%    \cs{@caption@uniquedo} (see above). \cs{@kernel@@caption} is a copy
%    of the kernel definition of \cs{@caption}: a package that has always
%    replaced \cs{caption} and \cs{@caption} (e.g., at
%    |\begin{document}|) can restore both and then use the hooks and
%    sockets above.
%
%    They are defined only if the kernel has the caption interface: after
%    a roll-back with \pkg{latexrelease} to an earlier date they are
%    undefined, and a roll-forward defines them again. The same holds for
%    \cs{@kernel@caption@reset}, which is defined together with
%    \cs{caption}. Therefore |\ifdefined\@kernel@caption| is the test for
%    the interface, for releases, roll-back, roll-forward and development
%    formats alike. There is no test whether \cs{caption} still is the
%    kernel's: a meaning test cannot tell wrappers from replacements.
%    \begin{macrocode}
%</2ekernel>
%<*2ekernel|latexrelease>
%<latexrelease>\IncludeInRelease{2026/11/01}%
%<latexrelease>                 {\@kernel@caption}{caption interface copies}%
\let\@kernel@caption\caption
\let\@kernel@caption@unique\@caption@uniquedo
\let\@kernel@@caption\@caption
%</2ekernel|latexrelease>
%<latexrelease>\EndIncludeInRelease
%<latexrelease>\IncludeInRelease{0000/00/00}%
%<latexrelease>                 {\@kernel@caption}{caption interface copies}%
%<latexrelease>\let\@kernel@caption\@undefined
%<latexrelease>\let\@kernel@caption@unique\@undefined
%<latexrelease>\let\@kernel@@caption\@undefined
%<latexrelease>\EndIncludeInRelease
%<*2ekernel>
%    \end{macrocode}
% \end{macro}
%
% \changes{v1.2k}{2026-10-06}{Hooks float/begin and float/end added}
% \paragraph{The hooks \texttt{float/begin} and \texttt{float/end}.}
%    Two hooks with two arguments each, for packages that need to
%    run code at the start and at the end of every float (for example
%    to set up float-type specific settings or a link target).
%    \begin{description}
%    \item[\texttt{float/begin}] is executed by \cs{@xfloat} (for
%      double-column floats in two-column mode by \cs{@xdblfloat}),
%      inside the box of the float and the group of the environment,
%      after \cs{@captype} has been set, \cs{@parboxrestore} and
%      \cs{@floatboxreset} have been executed and, for double-column
%      floats, \cs{hsize} and \cs{linewidth} have been set to
%      \cs{textwidth}.
%    \item[\texttt{float/end}] is executed by \cs{@endfloatbox}, inside
%      the box of the float, after the last paragraph has been ended
%      and before the box is closed.
%    \end{description}
%    The first argument is the float type (the value of \cs{@captype}),
%    the second argument is |*| for a double-column float in two-column
%    mode, and empty otherwise. The two hooks are a mirrored pair, so
%    code added to \texttt{float/end} is executed in reverse order.
%
%    The hooks are not related to the tagging sockets \texttt{float/begin}
%    and \texttt{float/end} declared in \texttt{lttagging.dtx}: sockets
%    and hooks are in different name spaces (the same holds for
%    \texttt{para/begin}). The sockets start and stop the structure
%    outside of the float box, the hooks are executed inside it.
%    \begin{macrocode}
%</2ekernel>
%<latexrelease>\IncludeInRelease{2026/11/01}%
%<latexrelease>                 {\@float@begin@hook}{Float hooks}%
%<*2ekernel|latexrelease>
\NewMirroredHookPairWithArguments{float/begin}{float/end}{2}
%    \end{macrocode}
%
% \begin{macro}{\@float@usehooks,\@float@begin@hook,\@float@end@hook,
%                \@float@hook@level}
% \changes{v1.2k}{2026-10-06}{Macros added}
%    \cs{@float@usehooks}\marg{type}\marg{star} executes the
%    \texttt{float/begin} hook and defines \cs{@float@end@hook} (locally,
%    inside the float box) to execute the \texttt{float/end} hook with the
%    same arguments. \cs{@endfloatbox} uses \cs{@float@end@hook}, which
%    is \cs{relax} outside of floats. Some packages use \cs{@endfloatbox}
%    for boxes that do not start with \cs{@xfloat} (for example
%    \pkg{float} for |[H]| floats), and such a box can be nested inside
%    a float, where the local definition is still visible. Therefore the
%    group level is recorded, and the end hook is executed only by the
%    \cs{@endfloatbox} at the same group level, that is, the one that
%    ends the box in which the begin hook was executed.
%    \cs{@float@begin@hook} is used by \cs{@xfloat}; \cs{@xdblfloat}
%    disables it and calls \cs{@float@usehooks} itself after the width
%    has been changed.
%    \begin{macrocode}
\def\@float@usehooks#1#2{%
  \edef\@float@hook@level{\the\currentgrouplevel}%
  \def\@float@end@hook{%
    \ifnum\currentgrouplevel=\@float@hook@level\relax
      \expandafter\@firstofone
    \else
      \expandafter\@gobble
    \fi
    {\UseHookWithArguments{float/end}{2}{#1}{#2}}}%
  \UseHookWithArguments{float/begin}{2}{#1}{#2}}
\def\@float@begin@hook#1{\@float@usehooks{#1}{}}
\let\@float@end@hook\relax
%</2ekernel|latexrelease>
%<latexrelease>\EndIncludeInRelease
%<latexrelease>\IncludeInRelease{0000/00/00}%
%<latexrelease>                 {\@float@begin@hook}{Float hooks}%
%<latexrelease>\let\@float@usehooks\@gobbletwo
%<latexrelease>\let\@float@begin@hook\@gobble
%<latexrelease>\let\@float@end@hook\relax
%<latexrelease>\EndIncludeInRelease
%<*2ekernel>
%    \end{macrocode}
% \end{macro}
%
% \begin{macro}{\@float}
% \begin{macro}{\@dblflset}
% \changes{v1.1a}{1994/10/31}{Macro added}
% \changes{v1.1g}{1994/12/10}{Macro removed temporarily}
% \changes{v1.1a}{1994/10/31}
%    {Major changes to parameter parsing, setting of local variables,
%    etc; two-column and one-column cases merged; space hacks moved}
% \changes{v1.1c}{1994/11/05}
%         {Add compatibility with old version of \cs{@xfloat}.}
% \changes{v1.1g}{1994/12/10}{Old version reinstated temporarily}
%
%    \begin{macrocode}
\def\@float#1{%
  \@ifnextchar[%
    {\@xfloat{#1}}%
    {\edef\reserved@a{\noexpand\@xfloat{#1}[\csname fps@#1\endcsname]}%
     \reserved@a}}
%    \end{macrocode}
%
% \end{macro}
% \end{macro}
%
%  \begin{macro}{\@dblfloat}
% \changes{v1.1a}{1994/10/31}
%     {Major changes since two-column and one-column cases merged}
% \changes{v1.1g}{1994/12/10}{Old version reinstated temporarily}
%
%    \begin{macrocode}
\def\@dblfloat{%
  \if@twocolumn\let\reserved@a\@dbflt\else\let\reserved@a\@float\fi
  \reserved@a}
%    \end{macrocode}
%  \end{macro}
%
%
%  \begin{macro}{\fps@dbl}
% \changes{v1.1a}{1994/10/31}{Macro added}
% \changes{v1.1g}{1994/12/10}{Macro removed temporarily}
%  Note that all double floats have default fps `tp'.
%  \end{macro}
%
%  \begin{macro}{\@setfps}
% \changes{v1.1a}{1994/10/31}{Macro added}
% \changes{v1.1c}{1994/11/05}
%         {Add compatibility with old version of \cs{@xfloat}.}
% \changes{v1.1g}{1994/12/10}{Macro removed temporarily}
%    This sets the fps, dealing with error conditions by adding
%    the default.
%
%  \end{macro}
%
%  \begin{macro}{\@xfloat}
% \changes{LaTeX2e}{1993/12/05}{Command changed}
% \changes{LaTeX2e}{1994/01/21}{Added missing percent characters.}
% \changes{v1.1a}{1994/10/31}
%     {Major changes, removing setting of local variables, space hacks
%     etc; two-column and one-column cases merged}
% \changes{v1.1c}{1994/11/05}
%         {Add compatibility with old version of \cs{@xfloat}: but the
%         arguments, provided at exorbitant cost, are now completely
%         ignored}
% \changes{v1.1f}{1994/11/21}
%     {Missing percents reinserted after 4, 8: these are not numbers.}
% \changes{v1.1g}{1994/12/10}{Old version reinstated temporarily}
% \changes{v1.1g}{1994/12/10}{Sanitization added temporarily}
%     The first part of this sets the count register that stores all
%     the information about the type and fps of the float.
%
%    We assume here that the default specifiers already contain no
%    active characters.
%
%    It may be better to store the defaults as numbers, rather than
%    symbol strings.
%
% \changes{v1.1p}{1996/10/24}{Added \cs{@nodocument} to trap
%                  floats in the preamble}
%    \begin{macrocode}
%</2ekernel>
%<latexrelease>\IncludeInRelease{2026/11/01}%
%<latexrelease>                 {\@xfloat}{Float hooks}%
%<*2ekernel|latexrelease>
\def\@xfloat #1[#2]{%
  \@nodocument
  \def \@captype {#1}%
   \def \@fps {#2}%
   \@onelevel@sanitize \@fps
   \def \reserved@b {!}%
   \ifx \reserved@b \@fps
     \@fpsadddefault
   \else
     \ifx \@fps \@empty
       \@fpsadddefault
     \fi
   \fi
   \ifhmode
     \@bsphack
     \@floatpenalty -\@Mii
   \else
     \@floatpenalty-\@Miii
   \fi
  \ifinner
     \@parmoderr\@floatpenalty\z@
  \else
    \@next\@currbox\@freelist
      {%
       \@tempcnta \sixt@@n
       \expandafter \@tfor \expandafter \reserved@a
         \expandafter :\expandafter =\@fps
         \do
%    \end{macrocode}
% \changes{v1.2b}{2015/01/11}{Check for valid option (latexrelease)}
% Start of changes, use a nested if structure, ending in an error.
%    \begin{macrocode}
          {%
           \if \reserved@a h%
             \ifodd \@tempcnta
             \else
               \advance \@tempcnta \@ne
             \fi
           \else\if \reserved@a t%
             \@setfpsbit \tw@
           \else\if \reserved@a b%
             \@setfpsbit 4%
           \else\if \reserved@a p%
             \@setfpsbit 8%
           \else\if \reserved@a !%
             \ifnum \@tempcnta>15
               \advance\@tempcnta -\sixt@@n\relax
             \fi
           \else
             \@latex@error{Unknown float option `\reserved@a'}%
             {Option `\reserved@a' ignored and `p' used.}%
             \@setfpsbit 8%
           \fi\fi\fi\fi\fi
           }%
%    \end{macrocode}
% End of changes
%    \begin{macrocode}
       \@tempcntb \csname ftype@\@captype \endcsname
       \multiply \@tempcntb \@xxxii
       \advance \@tempcnta \@tempcntb
       \global \count\@currbox \@tempcnta
       }%
    \@fltovf
  \fi
%    \end{macrocode}
%    The remainder sets up the box in which the float is typeset, and
%    the typesetting environment to be used.  It is essential to have
%    the extra box to avoid the unwanted space that would otherwise
%    often be put at the top of the float.
%
%    It ends with a hook; not sure how useful this is but it is needed
%    at present to deal with double-column floats.
% \task{CAR?}{Sort out hooks}
% \changes{v1.0a}{1994/03/07}
%     {(DPC) Extra group for color}
% \changes{v1.0c}{1994/03/14}
%     {(DPC) Use \cs{color@begingroup}}
% \changes{v1.0g}{1994/05/13}
%     {(DPC) Use \cs{normalcolor}}
% \changes{v1.1a}{1994/10/31}
%     {(DPC/CAR) Extra box added to remove color resetting from vmode}
% \changes{v1.1a}{1994/10/31}{Reset hook added}
% \changes{v1.1c}{1994/11/05}
%         {Use new \cs{color@hbox} concept.}
% \changes{v1.1f}{1994/11/21}
%         {Changed to \cs{color@vbox} so that large floats overflow
%          at the bottom}
% \changes{v1.1f}{1994/11/21}{Use \cs{@setnobreak}}
% \changes{v1.1f}{1994/11/21}{Added \cs{@setminipage}}
% \changes{v1.1f}{1994/11/21}{Added resetting of size and font}
% \changes{v1.1m}{1995/05/25}{(CAR) Resettings moved to hook}
%    \begin{macrocode}
  \global \setbox\@currbox
    \color@vbox
      \normalcolor
      \vbox \bgroup
        \hsize\columnwidth
        \@parboxrestore
        \@floatboxreset
%    \end{macrocode}
% \changes{v1.2k}{2026-10-06}{Execute the \texttt{float/begin} hook}
%    \begin{macrocode}
        \@float@begin@hook{#1}%
}%
%</2ekernel|latexrelease>
%<latexrelease>\EndIncludeInRelease
%<latexrelease>\IncludeInRelease{2015/01/01}%
%<latexrelease>                 {\@xfloat}{Check float options}%
%<latexrelease>\def\@xfloat #1[#2]{%
%<latexrelease>  \@nodocument
%<latexrelease>  \def \@captype {#1}%
%<latexrelease>   \def \@fps {#2}%
%<latexrelease>   \@onelevel@sanitize \@fps
%<latexrelease>   \def \reserved@b {!}%
%<latexrelease>   \ifx \reserved@b \@fps
%<latexrelease>     \@fpsadddefault
%<latexrelease>   \else
%<latexrelease>     \ifx \@fps \@empty
%<latexrelease>       \@fpsadddefault
%<latexrelease>     \fi
%<latexrelease>   \fi
%<latexrelease>   \ifhmode
%<latexrelease>     \@bsphack
%<latexrelease>     \@floatpenalty -\@Mii
%<latexrelease>   \else
%<latexrelease>     \@floatpenalty-\@Miii
%<latexrelease>   \fi
%<latexrelease>  \ifinner
%<latexrelease>     \@parmoderr\@floatpenalty\z@
%<latexrelease>  \else
%<latexrelease>    \@next\@currbox\@freelist
%<latexrelease>      {%
%<latexrelease>       \@tempcnta \sixt@@n
%<latexrelease>       \expandafter \@tfor \expandafter \reserved@a
%<latexrelease>         \expandafter :\expandafter =\@fps
%<latexrelease>         \do
%<latexrelease>          {%
%<latexrelease>           \if \reserved@a h%
%<latexrelease>             \ifodd \@tempcnta
%<latexrelease>             \else
%<latexrelease>               \advance \@tempcnta \@ne
%<latexrelease>             \fi
%<latexrelease>           \else\if \reserved@a t%
%<latexrelease>             \@setfpsbit \tw@
%<latexrelease>           \else\if \reserved@a b%
%<latexrelease>             \@setfpsbit 4%
%<latexrelease>           \else\if \reserved@a p%
%<latexrelease>             \@setfpsbit 8%
%<latexrelease>           \else\if \reserved@a !%
%<latexrelease>             \ifnum \@tempcnta>15
%<latexrelease>               \advance\@tempcnta -\sixt@@n\relax
%<latexrelease>             \fi
%<latexrelease>           \else
%<latexrelease>             \@latex@error{Unknown float option `\reserved@a'}%
%<latexrelease>             {Option `\reserved@a' ignored and `p' used.}%
%<latexrelease>             \@setfpsbit 8%
%<latexrelease>           \fi\fi\fi\fi\fi
%<latexrelease>           }%
%<latexrelease>       \@tempcntb \csname ftype@\@captype \endcsname
%<latexrelease>       \multiply \@tempcntb \@xxxii
%<latexrelease>       \advance \@tempcnta \@tempcntb
%<latexrelease>       \global \count\@currbox \@tempcnta
%<latexrelease>       }%
%<latexrelease>    \@fltovf
%<latexrelease>  \fi
%<latexrelease>  \global \setbox\@currbox
%<latexrelease>    \color@vbox
%<latexrelease>      \normalcolor
%<latexrelease>      \vbox \bgroup
%<latexrelease>        \hsize\columnwidth
%<latexrelease>        \@parboxrestore
%<latexrelease>        \@floatboxreset
%<latexrelease>}%
%<latexrelease>\EndIncludeInRelease
%<latexrelease>\IncludeInRelease{0000/00/00}%
%<latexrelease>                 {\@xfloat}{Check float options}%
%<latexrelease>\def\@xfloat #1[#2]{%
%<latexrelease>  \@nodocument
%<latexrelease>  \def \@captype {#1}%
%<latexrelease>   \def \@fps {#2}%
%<latexrelease>   \@onelevel@sanitize \@fps
%<latexrelease>   \def \reserved@b {!}%
%<latexrelease>   \ifx \reserved@b \@fps
%<latexrelease>     \@fpsadddefault
%<latexrelease>   \else
%<latexrelease>     \ifx \@fps \@empty
%<latexrelease>       \@fpsadddefault
%<latexrelease>     \fi
%<latexrelease>   \fi
%<latexrelease>   \ifhmode
%<latexrelease>     \@bsphack
%<latexrelease>     \@floatpenalty -\@Mii
%<latexrelease>   \else
%<latexrelease>     \@floatpenalty-\@Miii
%<latexrelease>   \fi
%<latexrelease>  \ifinner
%<latexrelease>     \@parmoderr\@floatpenalty\z@
%<latexrelease>  \else
%<latexrelease>    \@next\@currbox\@freelist
%<latexrelease>      {%
%<latexrelease>       \@tempcnta \sixt@@n
%<latexrelease>       \expandafter \@tfor \expandafter \reserved@a
%<latexrelease>         \expandafter :\expandafter =\@fps
%<latexrelease>         \do
%<latexrelease>          {%
%<latexrelease>           \if \reserved@a h%
%<latexrelease>             \ifodd \@tempcnta
%<latexrelease>             \else
%<latexrelease>               \advance \@tempcnta \@ne
%<latexrelease>             \fi
%<latexrelease>           \fi
%<latexrelease>           \if \reserved@a t%
%<latexrelease>             \@setfpsbit \tw@
%<latexrelease>           \fi
%<latexrelease>           \if \reserved@a b%
%<latexrelease>             \@setfpsbit 4%
%<latexrelease>           \fi
%<latexrelease>           \if \reserved@a p%
%<latexrelease>             \@setfpsbit 8%
%<latexrelease>           \fi
%<latexrelease>           \if \reserved@a !%
%<latexrelease>             \ifnum \@tempcnta>15
%<latexrelease>               \advance\@tempcnta -\sixt@@n\relax
%<latexrelease>             \fi
%<latexrelease>           \fi
%<latexrelease>           }%
%<latexrelease>       \@tempcntb \csname ftype@\@captype \endcsname
%<latexrelease>       \multiply \@tempcntb \@xxxii
%<latexrelease>       \advance \@tempcnta \@tempcntb
%<latexrelease>       \global \count\@currbox \@tempcnta
%<latexrelease>       }%
%<latexrelease>    \@fltovf
%<latexrelease>  \fi
%<latexrelease>  \global \setbox\@currbox
%<latexrelease>    \color@vbox
%<latexrelease>      \normalcolor
%<latexrelease>      \vbox \bgroup
%<latexrelease>        \hsize\columnwidth
%<latexrelease>        \@parboxrestore
%<latexrelease>        \@floatboxreset
%<latexrelease>}%
%<latexrelease>\EndIncludeInRelease
%<*2ekernel>
%    \end{macrocode}
%  \end{macro}
%
%  \begin{macro}{\@floatboxreset}
% \changes{v1.1a}{1994/10/31}{Macro added}
%
% The rational for allowing these normally global flags to be set
% locally here, via |\@parboxrestore|, was stated originally by
% Donald Arseneau and extended by Chris Rowley.
% It is because these flags are only set globally to
% true by section commands, and these should never appear within
% marginals or floats or, indeed, in any group; and they are only ever
% set globally to false when they are definitely true.
%
% If anyone is unhappy with this argument then both flags should be
% treated as in |\@setnobreak|; otherwise this command will be
% redundant.
% \changes{v1.1p}{1996/10/24}
%     {Added local settings of flags: dangerous!!}
%    \begin{macrocode}
\def \@floatboxreset {%
        \reset@font
        \normalsize
        \@setminipage
}
%    \end{macrocode}
%  \end{macro}
%
%  \begin{macro}{\@setnobreak}
% \changes{v1.1f}{1994/11/21}{Macro added}
% \changes{v1.1n}{1996/07/26}{remove unnecessary \cs{global} before
%                 \cs{@nobreak...}}
%    \begin{macrocode}
\def \@setnobreak{%
  \if@nobreak
    \let\outer@nobreak\@nobreaktrue
    \@nobreakfalse
  \fi
}
%    \end{macrocode}
%  \end{macro}
%
%  \begin{macro}{\@setminipage}
% \changes{v1.1f}{1994/11/21}{Macro added}
% \changes{v1.1n}{1996/07/26}{remove unnecessary \cs{global} before
%                 \cs{@minipage...}}
%    \begin{macrocode}
\def \@setminipage{%
  \@minipagetrue
  \everypar{\@minipagefalse\everypar{}}%
}
%    \end{macrocode}
%  \end{macro}
%
% \begin{macro}{\end@float}
% \changes{v1.0f}{1994/05/03}
%     {(CAR) Added \cs{@largefloatcheck}}
% \changes{v1.1n}{1995/10/25}{(CAR) unify code for double and
%    single versions}
%    \begin{macrocode}
\def\end@float{%
  \@endfloatbox
  \ifnum\@floatpenalty <\z@
%    \end{macrocode}
% We make sure that we never exceed |\textheight|, otherwise float
% will never get typeset (91/03/15 FMi).
%    \begin{macrocode}
    \@largefloatcheck
    \@cons\@currlist\@currbox
    \ifnum\@floatpenalty <-\@Mii
      \penalty -\@Miv
%    \end{macrocode}
% Saving and restoring |\prevdepth| added 26 May 87 to prevent extra
% vertical space when used in vertical mode.
%    \begin{macrocode}
      \@tempdima\prevdepth
      \vbox{}%
      \prevdepth\@tempdima
%    \end{macrocode}
%
%    \begin{macrocode}
      \penalty\@floatpenalty
%    \end{macrocode}
% \changes{LaTeX2.09}{1992/03/18}
%         {(RmS) changed \cs{@esphack} to \cs{@Esphack}}
%    \begin{macrocode}
    \else
      \vadjust{\penalty -\@Miv \vbox{}\penalty\@floatpenalty}\@Esphack
    \fi
  \fi
}
%    \end{macrocode}
% \end{macro}
%
% \begin{macro}{\end@dblfloat}
% \changes{v1.0f}{1994/05/03}{\cs{@largefloatcheck} added}
% \changes{v1.1n}{1995/10/25}{(CAR) unify code for double and
%    single versions}
% \changes{v1.2b}{2015/01/11}{float order in 2-column (latexrelease)}
%    \begin{macrocode}
%</2ekernel>
%<latexrelease>\IncludeInRelease{2015/01/01}%
%<latexrelease>                 {\end@dblfloat}{float order in 2-column}%
%<*2ekernel|latexrelease>
\def\end@dblfloat{%
  \if@twocolumn
    \@endfloatbox
    \ifnum\@floatpenalty <\z@
      \@largefloatcheck
%    \end{macrocode}
%
%  Force the depth of two column float boxes.
%    \begin{macrocode}
      \global\dp\@currbox1sp %
%    \end{macrocode}
%    What follows is essentially |\end@float| without a starting
%    |\@endfloatbox|.
% \changes{v1.2b}{2000/09/24}{FMi: use output routine to
%    defer float}
% \changes{v1.2b}{2014/04/27}{Inline the code to allow some
%   coexistence with packages that hook into \cs{end@float} and do not
%   know about the algorithm change}
%    \begin{macrocode}
      \@cons\@currlist\@currbox
      \ifnum\@floatpenalty <-\@Mii
        \penalty -\@Miv
        \@tempdima\prevdepth
        \vbox{}%
        \prevdepth\@tempdima
        \penalty\@floatpenalty
      \else
        \vadjust{\penalty -\@Miv \vbox{}\penalty\@floatpenalty}\@Esphack
      \fi
%    \end{macrocode}
% \changes{v1.2b}{2014/06/10}{missing \cs{fi} added}
%    \begin{macrocode}
    \fi
  \else
    \end@float
  \fi
}%
%</2ekernel|latexrelease>
%<latexrelease>\EndIncludeInRelease
%<latexrelease>\IncludeInRelease{0000/00/00}%
%<latexrelease>                 {\end@dblfloat}{float order in 2-column}%
%<latexrelease>\def\end@dblfloat{%
%<latexrelease>\if@twocolumn
%<latexrelease>  \@endfloatbox
%<latexrelease>  \ifnum\@floatpenalty <\z@
%    \end{macrocode}
% We make sure that we never exceed |\textheight|, otherwise float
% will never get typeset (91/03/15 FMi).
%    \begin{macrocode}
%<latexrelease>    \@largefloatcheck
%<latexrelease>    \@cons\@dbldeferlist\@currbox
%<latexrelease>  \fi
%    \end{macrocode}
% RmS 92/03/18 changed |\@esphack| to |\@Esphack|.
%    \begin{macrocode}
%<latexrelease>    \ifnum \@floatpenalty =-\@Mii \@Esphack\fi
%<latexrelease>\else
%<latexrelease>  \end@float
%<latexrelease>\fi
%<latexrelease>}%
%<latexrelease>\EndIncludeInRelease
%<*2ekernel>
%    \end{macrocode}
% \end{macro}
%
% \begin{macro}{\@endfloatbox}
% \changes{v1.1n}{1995/10/25}{(CAR) macro added: to unify code for
%    double and single versions}
%    This macro is not intended to be a hook; it is designed to help
%    maintain the integrity of this code, which is used twice and, as
%    can be seen, is subject to frequent changes.
%    \begin{macrocode}
%</2ekernel>
%<latexrelease>\IncludeInRelease{2026/11/01}%
%<latexrelease>                 {\@endfloatbox}{Float hooks}%
%<*2ekernel|latexrelease>
\def \@endfloatbox{%
      \par
%    \end{macrocode}
% \changes{v1.2k}{2026-10-06}{Execute the \texttt{float/end} hook}
%    \begin{macrocode}
      \@float@end@hook
      \vskip\z@skip      %% \par\vskip\z@ added 15 Dec 87
%    \end{macrocode}
% \changes{v1.0a}
%     {1994/03/07}{(DPC) Extra group for color}
% \changes{v1.0c}{1994/03/14}
%     {(DPC) Use \cs{color@endgroup}}
% \changes{v1.0h}{1994/05/20}{Restore outer value of @nobreak switch.}
% \changes{v1.1a}{1994/10/31}
%     {(DPC/CAR) Extra box added to remove color resetting from vmode}
% \changes{v1.1c}{1994/11/05}
%         {Use new \cs{color@hbox} concept.}
% \changes{v1.1f}{1994/11/21}{Corrected position of \cs{outer@nobreak}}
% \changes{v1.1f}{1994/11/21}{Added reset of minipage flag}
% \changes{v1.1n}{1996/07/26}{remove unnecessary \cs{global} before
%                 \cs{@minipage...}}
%    \begin{macrocode}
      \@minipagefalse
      \outer@nobreak
    \egroup                  %% end of vbox
  \color@endbox
}
%</2ekernel|latexrelease>
%<latexrelease>\EndIncludeInRelease
%<latexrelease>\IncludeInRelease{0000/00/00}%
%<latexrelease>                 {\@endfloatbox}{Float hooks}%
%<latexrelease>\def \@endfloatbox{%
%<latexrelease>      \par\vskip\z@skip
%<latexrelease>      \@minipagefalse
%<latexrelease>      \outer@nobreak
%<latexrelease>    \egroup
%<latexrelease>  \color@endbox
%<latexrelease>}
%<latexrelease>\EndIncludeInRelease
%<*2ekernel>
%    \end{macrocode}
% \end{macro}
%
% \begin{macro}{\outer@nobreak}
% \changes{v1.0h}{1994/05/20}{Macro added: default is to do nothing.}
%    \begin{macrocode}
\let\outer@nobreak\@empty
%    \end{macrocode}
% \end{macro}
%
%
%  \begin{macro}{\@largefloatcheck}
% \changes{v1.0e}{1994/04/25}{Command added}
%
%    This calculates by how much a float is oversize for the page and
%    prints this in a warning message.
%
%    \begin{macrocode}
\def \@largefloatcheck{%
  \ifdim \ht\@currbox>\textheight
    \@tempdima -\textheight
    \advance \@tempdima \ht\@currbox
%    \end{macrocode}
% \changes{v1.0e}{1994/04/25}{Changed warning message to give more
% info}
%    \begin{macrocode}
    \@latex@warning {Float too large for page by \the\@tempdima}%
    \ht\@currbox \textheight
  \fi
}
%    \end{macrocode}
%  \end{macro}
%
%  \begin{macro}{\@dbflt}
%  \begin{macro}{\@xdblfloat}
% \changes{v1.1a}{1994/10/31}
%     {Macros removed: \cs{@dbflt}, \cs{@xdblfloat}}
% \changes{v1.1g}{1994/12/10}{Macros reinserted temporarily}
%
%
% \changes{v1.2k}{2026-10-06}{Execute the \texttt{float/begin} hook
%    after the width has been set}
%    \begin{macrocode}
\def\@dbflt#1{\@ifnextchar[{\@xdblfloat{#1}}{\@xdblfloat{#1}[tp]}}
%</2ekernel>
%<latexrelease>\IncludeInRelease{2026/11/01}%
%<latexrelease>                 {\@xdblfloat}{Float hooks}%
%<*2ekernel|latexrelease>
\def\@xdblfloat#1[#2]{%
  \let\@float@begin@hook\@gobble
  \@xfloat{#1}[#2]\hsize\textwidth\linewidth\textwidth
  \@float@usehooks{#1}{*}}
%</2ekernel|latexrelease>
%<latexrelease>\EndIncludeInRelease
%<latexrelease>\IncludeInRelease{0000/00/00}%
%<latexrelease>                 {\@xdblfloat}{Float hooks}%
%<latexrelease>\def\@xdblfloat#1[#2]{%
%<latexrelease>  \@xfloat{#1}[#2]\hsize\textwidth\linewidth\textwidth}
%<latexrelease>\EndIncludeInRelease
%<*2ekernel>
%    \end{macrocode}
%  \end{macro}
%  \end{macro}
%
%
% Moved to ltoutput 93/12/16
%    \begin{macrocode}
%\newcount\c@topnumber
%\newcount\c@dbltopnumber
%\newcount\c@bottomnumber
%\newcount\c@totalnumber
%    \end{macrocode}
%
%  \begin{macro}{\@floatplacement}
% An analysis of |\@floatplacement|:
%
% This should be called whenever |\@colht| has been set.
%    \begin{macrocode}
\def\@floatplacement{\global\@topnum\c@topnumber
    % Textpage bit, global:
   \global\@toproom \topfraction\@colht
   \global\@botnum  \c@bottomnumber
   \global\@botroom \bottomfraction\@colht
   \global\@colnum  \c@totalnumber
    % Floatpage bit, local:
   \@fpmin   \floatpagefraction\@colht}
%</2ekernel>
%    \end{macrocode}
% \end{macro}
%
%  \begin{macro}{\@dblfloatplacement}
% \changes{LaTeX2e}{1993/12/05}{Command changed}
%
%     This should be called only within a group.  Now changed to
%     provide extra checks in |\@addtodblcol|, needed when processing a
%     BANG float.
%
% \changes{v1.2b}{2015/01/11}{float order in 2-column (latexrelease)}
%    \begin{macrocode}
%<latexrelease>\IncludeInRelease{2015/01/01}%
%<latexrelease>          {\@dblfloatplacement}{float order in 2-column}%
%<*2ekernel|latexrelease>
%    \end{macrocode}
% When making two column float area, look for floats with 1sp
% depth.
%    \begin{macrocode}
\def\@dblfloatplacement{\global\@dbltopnum\c@dbltopnumber
   \global\@dbltoproom \dbltopfraction\@colht
   \@textmin \@colht
   \advance \@textmin -\@dbltoproom
   \@fpmin \dblfloatpagefraction\textheight
   \@fptop \@dblfptop
   \@fpsep \@dblfpsep
   \@fpbot \@dblfpbot
%    \end{macrocode}
%    |\f@depth| is used in |\@testwrongwidth| to look for either
%   column or dbl-column floats. A value of |1sp| signals the
%   latter. Because of this setting here, |\@dblfloatplacment| needs to
%   be called inside a group which is a questionable design.
%    \begin{macrocode}
   \def\f@depth{1sp}}%
%    \end{macrocode}
%
%    \begin{macrocode}
%</2ekernel|latexrelease>
%<latexrelease>\EndIncludeInRelease
%<latexrelease>\IncludeInRelease{0000/00/00}%
%<latexrelease>          {\@dblfloatplacement}{float order in 2-column}%
%<latexrelease>\def \@dblfloatplacement {%
%    \end{macrocode}
%    Textpage bit: global, but need not be.
%    \begin{macrocode}
%<latexrelease>  \global \@dbltopnum \c@dbltopnumber
%<latexrelease>  \global \@dbltoproom \dbltopfraction\@colht
%    \end{macrocode}
%   This new bit uses |\@textmin| to locally store the amount of extra
%   room in the column.
%    \begin{macrocode}
%<latexrelease>  \@textmin \@colht
%<latexrelease>  \advance \@textmin -\@dbltoproom
%    \end{macrocode}
%    Floatpage bit: must be local.
%    \begin{macrocode}
%<latexrelease>  \@fpmin \dblfloatpagefraction\textheight
%<latexrelease>  \@fptop \@dblfptop
%<latexrelease>  \@fpsep \@dblfpsep
%<latexrelease>  \@fpbot \@dblfpbot
%<latexrelease>}%
%<latexrelease>\EndIncludeInRelease
%<*2ekernel>
%    \end{macrocode}
%  \end{macro}
%
%
% \begin{oldcomments}
%   MARGINAL NOTES:
%
%   Marginal notes use the same mechanism as floats to communicate
%   with the \output routine.  Marginal notes are distinguished from
%   floats by having a negative placement specification.  The command
%   \marginpar [LTEXT]{RTEXT} generates a marginal note in a parbox,
%   using LTEXT if it's on the left and RTEXT if it's on the right.
%   (Default is RTEXT = LTEXT.)  It uses the following parameters.
%
%   \marginparwidth : Width of marginal notes.
%   \marginparsep   : Distance between marginal note and text.
%        the page layout to determine how to move the marginal
%        note into the margin.   E.g., \@leftmarginskip ==
%        \hskip -\marginparwidth \hskip -\marginparsep .
%   \marginparpush  :  Minimum vertical separation between \marginpar's
%
%  Marginal notes are normally put on the outside of the page
%  if @mparswitch = true, and on the right if @mparswitch = false.
%  The command \reversemarginpar reverses the side where they
%  are put.  \normalmarginpar undoes \reversemarginpar.
%  These commands have no effect for two-column output.
%
%  SURPRISE: if two marginal notes appear on the same line of
%  text, then the second one could appear on the next page, in
%  a funny position.
%
%
%  \marginpar [LTEXT]{RTEXT} ==
%   BEGIN
%     if hmode then \@bsphack
%                   \@floatpenalty := -10002
%              else \@floatpenalty := -10003
%     fi
%     if inner
%       then LaTeX Error: 'Not in outer paragraph mode.'
%            \@floatpenalty := 0
%       else if \@freelist has two elements:
%              then get \@marbox, \@currbox  from \@freelist
%                   \count\@marbox :=G -1
%              else \@floatpenalty := 0
%                   LaTeX Error: 'Too many unprocessed floats'
%                   \@currbox, \@marbox := \@tempboxa    %%use \def
%            fi
%     fi
%     if optional argument
%       then %% \@xmpar ==
%            \@savemarbox\@marbox{LTEXT}
%            \@savemarbox\@currbox{RTEXT}
%       else %% \@ympar ==
%            \@savemarbox\@marbox{RTEXT}
%            \box\@currbox :=G \box\@marbox
%    fi
%    \@xympar
%   END
%
% \reversemarginpar == BEGIN \@mparbottom   :=G 0
%                            @reversemargin :=G true
%                      END
%
% \normalmarginpar  == BEGIN \@mparbottom   :=G 0
%                            @reversemargin :=G false
%                      END
%
% \end{oldcomments}
%
% \begin{macro}{\marginpar}
%    \begin{macrocode}
\def\marginpar{%
  \ifhmode
    \@bsphack
    \@floatpenalty -\@Mii
  \else
    \@floatpenalty-\@Miii
  \fi
  \ifinner
    \@parmoderr
    \@floatpenalty\z@
  \else
    \@next\@currbox\@freelist{}{}%
    \@next\@marbox\@freelist{\global\count\@marbox\m@ne}%
       {\@floatpenalty\z@
        \@fltovf\def\@currbox{\@tempboxa}\def\@marbox{\@tempboxa}}%
  \fi
  \@ifnextchar [\@xmpar\@ympar}
%    \end{macrocode}
% \end{macro}
%
% \begin{macro}{\@xmpar}
%    \begin{macrocode}
\long\def\@xmpar[#1]#2{%
  \@savemarbox\@marbox{#1}%
  \@savemarbox\@currbox{#2}%
  \@xympar}
%    \end{macrocode}
% \end{macro}
%
% \begin{macro}{\@ympar}
%    \begin{macrocode}
\long\def\@ympar#1{%
  \@savemarbox\@marbox{#1}%
  \global\setbox\@currbox\copy\@marbox
  \@xympar}
%    \end{macrocode}
% \end{macro}
%
% \begin{macro}{\@savemarbox}
% \changes{v1.0b}{1994/03/12}
%     {(DPC) Extra group for color}
% \changes{v1.0c}{1994/03/14}
%     {(DPC) Use \cs{color@begingroup}}
% \changes{v1.0d}{1994/04/18}
%     {(DPC) Remove Color support}
% \changes{v1.1a}{1994/10/31}
%     {(DPC/CAR) Extra box added for color}
% \changes{v1.1c}{1994/11/05}
%         {Use new \cs{color@hbox} concept.}
% \changes{v1.1f}{1994/11/21}{Changed to \cs{color@vbox} }
% \changes{v1.1f}{1994/11/21}{Use \cs{@setnobreak} etc}
% \changes{v1.1f}{1994/11/21}{Added \cs{@setminipage} etc}
% \changes{v1.1f}{1994/11/21}{Added resetting of size and font}
% \changes{v1.1m}{1995/05/25}{(CAR) Resettings moved to hook}
% \changes{v1.1n}{1996/07/26}{remove unnecessary \cs{global} before
%                 \cs{@minipage...}}
% \changes{v1.2e}{2021/02/03}{Explicitly end with \cs{par} (gh/489)}
%    \begin{macrocode}
%</2ekernel>
%<*2ekernel|latexrelease>
%<latexrelease>\IncludeInRelease{2021/06/01}%
%<latexrelease>                 {\@savemarbox}{Explicit par for marginpar}%
\long\def \@savemarbox #1#2{%
  \global\setbox #1%
    \color@vbox
      \vtop{%
        \hsize\marginparwidth
        \@parboxrestore
        \@marginparreset
        #2\par
        \@minipagefalse
        \outer@nobreak
        }%
    \color@endbox
}
%</2ekernel|latexrelease>
%<latexrelease>\EndIncludeInRelease
%    \end{macrocode}
%
%    \begin{macrocode}
%<latexrelease>\IncludeInRelease{0000/00/00}%
%<latexrelease>                 {\@savemarbox}{Explicit par for marginpar}%
%<latexrelease>
%<latexrelease>\long\def \@savemarbox #1#2{%
%<latexrelease>  \global\setbox #1%
%<latexrelease>    \color@vbox
%<latexrelease>      \vtop{%
%<latexrelease>        \hsize\marginparwidth
%<latexrelease>        \@parboxrestore
%<latexrelease>        \@marginparreset
%<latexrelease>        #2%
%<latexrelease>        \@minipagefalse
%<latexrelease>        \outer@nobreak
%<latexrelease>        }%
%<latexrelease>    \color@endbox
%<latexrelease>}
%<latexrelease>\EndIncludeInRelease
%<*2ekernel>
%    \end{macrocode}
% \end{macro}
%
%  \begin{macro}{\@marginparreset}
% \changes{v1.1f}{1994/11/21}{Macro added}
%
% The rational for allowing these normally global flags to be set
% locally here, via |\@parboxrestore| was stated originally by
% Donald Arseneau and extended by Chris Rowley.
% It is because these flags are only set globally to
% true by section commands, and these should never appear within
% marginals or floats or, indeed, in any group; and they are only ever
% set globally to false when they are definitely true.
%
% If anyone is unhappy with this argument then both flags should be
% treated as in |\@setnobreak|; otherwise this command will be
% redundant.
% \changes{v1.1p}{1996/10/24}
%     {Added local settings of flags: dangerous!!}
%    \begin{macrocode}
\def \@marginparreset {%
        \reset@font
        \normalsize
%        \let\if@nobreak\iffalse
%        \let\if@noskipsec\iffalse
%        \@setnobreak
        \@setminipage
}
%    \end{macrocode}
%  \end{macro}
%
% \begin{macro}{\@xympar}
%    \begin{macrocode}
%    \end{macrocode}
% \changes{LaTeX2.09}{1992/03/18}
%     {(RmS) added \cs{global}\cs{@ignorefalse}}
% \changes{v1.0b}{1994/03/12}
%     {(DPC) Extra bgroup for color}
% \changes{1.0c}{1994/03/14}
%     {(DPC) Use \cs{color@begingroup}}
% \changes{v1.1a}{1994/10/31}
%     {(DPC/CAR) Extra box added since needed for floats}
% \changes{v1.1c}{1994/11/05}
%         {Use new \cs{color@hbox} concept.}
% \changes{v1.1f}{1994/11/21}{Changed to \cs{color@vbox} }
%     Setting the box here is done only because the code
%     uses \cs{end@float}; it will be empty and gets discarded.
% \changes{v1.1o}{1996/08/02}{Remove \cs{global} before \cs{@ignore...}}
%    \begin{macrocode}
\def \@xympar{%
  \ifnum\@floatpenalty <\z@\@cons\@currlist\@marbox\fi
  \setbox\@tempboxa
    \color@vbox
      \vbox \bgroup
  \end@float
  \@ignorefalse
  \@esphack
}
%    \end{macrocode}
% \end{macro}
%
% \begin{macro}{\reversemarginpar}
% \begin{macro}{\normalmarginpar}
%    \begin{macrocode}
\def\reversemarginpar{\global\@mparbottom\z@ \@reversemargintrue}
\def\normalmarginpar{\global\@mparbottom\z@ \@reversemarginfalse}
%    \end{macrocode}
% \end{macro}
% \end{macro}
%
%    \begin{macrocode}
\message{footnotes,}
%    \end{macrocode}
%
% \subsection{Footnotes}
%
% \begin{oldcomments}
%
%   \footnote{NOTE}       : User command to insert a footnote.
%
%   \footnote[NUM]{NOTE}: User command to insert a footnote numbered
%                       NUM, where NUM is a number -- 1, 2,
%                       etc.  For example, if footnotes are numbered
%                       *, **, etc. within pages, then \footnote[2]{...}
%                       produces footnote '**'.  This command does not
%                       step the footnote counter.
%
%   \footnotemark[NUM] : Command to produce just the footnote mark in
%                        the text, but no footnote.  With no argument,
%                        it steps the footnote counter before generating
%                        the mark.
%
%   \footnotetext[NUM]{TEXT} : Command to produce the footnote but
%                              no mark.  \footnote is equivalent to
%                              \footnotemark \footnotetext .
%
%   As in PLAIN, footnotes use \insert\footins, and the following
%   parameters:
%
%   \footnotesize   : Size-changing command for footnotes.
%
%   \footnotesep    : The height of a strut placed at the beginning of
%                     every footnote.
%   \skip\footins   : Space between main text and footnotes.  The rule
%                     separating footnotes from text occurs in this
%                     space. This space lies above the strut of height
%                     \footnotesep which is at the beginning of the
%                     first footnote.
%   \footnoterule   : Macro to draw the rule separating footnotes from
%                     text. It is executed right after a \vspace of
%                     \skip\footins. It should take zero vertical
%                     space--i.e., it should to a negative skip to
%                     compensate for any positive space it occupies.
%                     (See PLAIN.TEX.)
%
%   \interfootnotelinepenalty : Interline penalty for footnotes.
%
%   \thefootnote : In usual LaTeX style, produces the footnote number.
%                  If footnotes are to be numbered within pages, then
%                  the document style file must include an \@addtoreset
%                  command to cause the footnote counter to be reset
%                  when the page counter is stepped.  This is not a good
%                  idea, though, because the counter will not always be
%                  reset in time to ensure that the first footnote on a
%                  page is footnote number one.
%
%   \@thefnmark : Holds the current footnote's mark--e.g., \dag or '1'
%                 or 'a'.
%
%   \@mpfnnumber  : A macro that generates the numbers for \footnote
%                  and \footnotemark commands. It == \thefootnote
%                  outside a minipage environment, but can be
%                  changed inside to generate numbers for
%                  \footnote's.
%
%   \@makefnmark : A macro to generate the footnote marker from
%                 \@thefnmark The default definition was
%                 \hbox{$^\@thefnmark$}.
%
%                 This is now replaced by
%                 \textsuperscript{\@thefnmark}
%
%   \@makefntext{NOTE} :
%        Must produce the actual footnote, using \@thefnmark as the mark
%        of the footnote and NOTE as the text.  It is called when
%        effectively  inside a \parbox, with \hsize = \columnwidth.
%          For example, it might be as simple as
%               $^{\@thefnmark}$  NOTE
%
% In a minipage environment, \footnote and \footnotetext are redefined
% so that
%    (a) they use the counter mpfootnote
%    (b) the footnotes they produce go at the bottom of the minipage.
% The switch is accomplished by letting \@mpfn == footnote or mpfootnote
% and \thempfn == \thefootnote or \thempfootnote, and by redefining
% \@footnotetext to be \@mpfootnotetext in the minipage.
%
% \footnote{NOTE}  ==
%  BEGIN
%    \stepcounter{\@mpfn}
%    begingroup
%       \protect == \noexpand
%       \@thefnmark :=G eval (\thempfn)
%    endgroup
%    \@footnotemark
%    \@footnotetext{NOTE}
%  END
%
% \footnote[NUM]{NOTE} ==
%  BEGIN
%    begingroup
%       \protect == \noexpand
%       counter \@mpfn :=L NUM
%       \@thefnmark :=G eval (\thempfn)
%    endgroup
%    \@footnotemark
%    \@footnotetext{NOTE}
%  END
%
% \footnotemark      ==
%  BEGIN \stepcounter{footnote}
%        begingroup
%           \protect == \noexpand
%           \@thefnmark :=G eval(\thefootnote)
%        endgroup
%        \@footnotemark
%  END
%
% \footnotemark[NUM] ==
%   BEGIN
%       begingroup
%         footnote counter :=L NUM
%         \protect == \noexpand
%        \@thefnmark :=G eval(\thefootnote)
%       endgroup
%       \@footnotemark
%   END
%
% \@footnotemark ==
%   BEGIN
%    \leavevmode
%    IF hmode THEN \@x@sf := \the\spacefactor FI
%    \@makefnmark          % put number in main text
%    IF hmode THEN \spacefactor := \@x@sf FI
%   END
%
% \footnotetext      ==
%    BEGIN begingroup \protect == \noexpand
%                     \@thefnmark :=G eval (\thempfn)
%          endgroup
%          \@footnotetext
%    END
%
% \footnotetext[NUM] ==
%    BEGIN begingroup  counter \@mpfn :=L NUM
%                      \protect == \noexpand
%                      \@thefnmark :=G eval (\thempfn)
%          endgroup
%          \@footnotetext
%    END
%
% \end{oldcomments}
%
%
% \changes{v1.1l}{1995/05/24}{Moved definition of \cs{footins}
%                   and \cs{footnoterule} from ltplain.}
%
% \begin{macro}{\footins}
% \LaTeX\ does use the same insert for footnotes as PLAIN.
%    \begin{macrocode}
\newinsert\footins
%    \end{macrocode}
%
% \LaTeX\ leaves these initializations for the |\footins| insert.
%
%    \begin{macrocode}
\skip\footins=\bigskipamount % space added when footnote is present
\count\footins=1000 % footnote magnification factor (1 to 1)
\dimen\footins=8in % maximum footnotes per page
%    \end{macrocode}
% \end{macro}
%
%
% \begin{macro}{\footnoterule}
% \LaTeX\ keeps PLAIN \TeX's |\footnoterule| as the default.
%
%    \begin{macrocode}
\def\footnoterule{\kern-3\p@
  \hrule \@width 2in \kern 2.6\p@} % the \hrule is .4pt high
%    \end{macrocode}
% \end{macro}
%
% \begin{macro}{\thefootnote}
% \changes{v1.1i}{1995/05/16}{Streamlined parts of code.}
%    \begin{macrocode}
\@definecounter{footnote}
\def\thefootnote{\@arabic\c@footnote}
%    \end{macrocode}
% \end{macro}
%
% \begin{macro}{\thempfootnote}
% \changes{v1.1j}{1995/05/18}{Added \cs{itshape}.}
% \changes{v1.1v}{2002/10/01}{Use braces around \cs{itshape}
%    to keep font change local (pr/3460).}
%    The default display for the footnote counter in minipages is to
%    use italic letters. We use |\itshape| not |\textit| as the latter
%    would add an italic correction.
%    \begin{macrocode}
\@definecounter{mpfootnote}
\def\thempfootnote{{\itshape\@alph\c@mpfootnote}}
%    \end{macrocode}
% \end{macro}
%
% \begin{macro}{\@makefnmark}
% \changes{v1.1i}{1995/05/16}{Now use \cs{textsuperscript}.}
% \changes{v1.1j}{1995/05/18}{Added \cs{normalfont}.}
% \changes{v1.1k}{1995/05/20}{Moved \cs{normalfont} to
%                             \cs{textsuperscript}}
% \changes{v1.1k}{1995/05/20}{Moved \cs{normalfont} back
%                        and use \cs{@textsuperscript}}
%    Default definition.
%    \begin{macrocode}
%\def\@makefnmark{\hbox{$^{\@thefnmark}\m@th$}}
\def\@makefnmark{\hbox{\@textsuperscript{\normalfont\@thefnmark}}}
%    \end{macrocode}
% \end{macro}
%
%  \begin{macro}{\textsuperscript}
% \changes{v1.1i}{1995/05/16}{Command added./pr1503}
% \changes{v1.1k}{1995/05/20}{Use \cs{normalfont}.}
% \changes{v1.1l}{1995/05/24}{Use \cs{@textsuperscript}}
%    This command provides superscript characters in the current text
%    font. It's implementation might change!!!
%    \begin{macrocode}
\DeclareRobustCommand*\textsuperscript[1]{%
  \@textsuperscript{\selectfont#1}}
%    \end{macrocode}
%  \end{macro}
%
%  \begin{macro}{\@textsuperscript}
% \changes{v1.1l}{1995/05/24}{Command added.}
% \changes{v1.1n}{1995/12/05}{Use \cs{ensuremath} for latex/1984.}
% \changes{v1.1m}{1995/12/07}
%      {Move \cs{m@th} out of the \cs{ensuremath} for latex/1984.}
% \changes{v1.2h}{2026-02-15}
%      {Drop use of fake math}
% \changes{v1.2j}{2026-10-01}{Move \cs{check@mathfonts} inside group}
%  \begin{macro}{\textsuperscript@offset}
% \changes{v1.2h}{2026-02-15}{Command added}
%  \begin{macro}{\textsuperscript@space}
% \changes{v1.2h}{2026-02-15}{Command added}
%    This command should not be used directly, but may be used to define
%   other commands |\textsuperscript|, |\@makefnmark|. |#1| should
%   always start with a font selection command, to activate the font
%   size switch. Doing everything in text mode using the classical dimensions
%   means using the settings from math mode: to allow the user to deviate
%   from those, we use a one-level indirection. Notice that the script space
%   is applied without using a kern to keep the tracing output as similar as
%   possible to the previous math-based approach. The rather involved
%   default for |\textsuperscript@offset| arises so we get the same result
%   as from the math mode predecessor. The one for |\textsuperscript@space|
%   is needed as in math mode, |\scriptspace| is not dropped at a line break
%   (as it's a box adjustment). Lua\TeX{} has rather easier to understand
%   names for the various parameters, and these will be correct whether
%   classical or OpenType fonts are in use for math, so we set up with
%   these if available.
%    \begin{macrocode}
%</2ekernel>
%<latexrelease>\IncludeInRelease{2026/06/01}%
%<latexrelease>                 {\@textsuperscript}{real text superscript}%
%<*2ekernel|latexrelease>
\protected\def\@textsuperscript#1%
  {%
    \leavevmode
    \begingroup
      \check@mathfonts
      \sbox\z@{\fontsize\sf@size\sf@size#1}%
      \raise\dimexpr\textsuperscript@offset\relax\box\z@
      \textsuperscript@space
    \endgroup
  }
\ifdefined\directlua
  \def\textsuperscript@offset{%
    \ifdim\Umathsupshiftup\textstyle<\dimexpr\dp\z@+\Umathsupbottommin\textstyle\relax
      \dp\z@+\Umathsupbottommin\textstyle
    \else
      \Umathsupshiftup\textstyle
    \fi
  }
\else
  \def\textsuperscript@offset{%
    \ifdim\fontdimen14\textfont\tw@<\dimexpr\dp\z@+0.25\fontdimen5\textfont\tw@\relax
      \dp\z@+0.25\fontdimen5\textfont\tw@
    \else
      \fontdimen14\textfont\tw@
    \fi
  }
\fi
\def\textsuperscript@space{\nobreak\hskip\scriptspace\kern\z@}
%</2ekernel|latexrelease>
%<latexrelease>\EndIncludeInRelease
%    \end{macrocode}
%
%    \begin{macrocode}
%<latexrelease>\IncludeInRelease{2020/10/01}%
%<latexrelease>                 {\@textsuperscript}{real text superscript}%
%<latexrelease>\def\@textsuperscript#1{%
%<latexrelease>  {\m@th\ensuremath{^{\mbox{\fontsize\sf@size\sf@size#1}}}}}
%<latexrelease>\EndIncludeInRelease
%    \end{macrocode}
%
%    \begin{macrocode}
%<latexrelease>\IncludeInRelease{0000/00/00}%
%<latexrelease>                 {\@textsuperscript}{real text superscript}%
%<latexrelease>
%<latexrelease>\def\@textsuperscript#1{%
%<latexrelease>  {\m@th\ensuremath{^{\mbox{\fontsize\sf@size\z@#1}}}}}
%<latexrelease>\EndIncludeInRelease
%<*2ekernel>
%    \end{macrocode}
%  \end{macro}
%  \end{macro}
%  \end{macro}
%
%
%
%
%
%
%  \begin{macro}{\textsubscript}
% \changes{v1.2a}{2014/12/30}{Command added (latexrelease)}
% \changes{v1.2d}{2020/04/09}{Set non-zero baseline (gh/249)}
%    \begin{macrocode}
%</2ekernel>
%<latexrelease>\IncludeInRelease{2015/01/01}%
%<latexrelease>                 {\textsubscript}{\textsubscript}%
%<*2ekernel|latexrelease>
%    \end{macrocode}
%
%    \begin{macrocode}
\DeclareRobustCommand*\textsubscript[1]{%
  \@textsubscript{\selectfont#1}}%
%    \end{macrocode}
%
%    \begin{macrocode}
%</2ekernel|latexrelease>
%<latexrelease>\EndIncludeInRelease
%<latexrelease>\IncludeInRelease{0000/00/00}%
%<latexrelease>                 {\textsubscript}{\textsubscript}%
%<latexrelease>\let\textsubscript\@undefined
%<latexrelease>\EndIncludeInRelease
%<*2ekernel>
%    \end{macrocode}
%  \end{macro}
%
%
%
%
%  \begin{macro}{\@textsubscript}
% \changes{v1.2a}{2014/12/30}{Command added (latexrelease)}
% \changes{v1.2d}{2020/04/09}{Set non-zero baseline (gh/249)}
% \changes{v1.2h}{2026-02-15}
%      {Drop use of fake math}
% \changes{v1.2j}{2026-10-01}{Move \cs{check@mathfonts} inside group}
%  \begin{macro}{\textsubscript@offset}
% \changes{v1.2h}{2026-02-15}{Command added}
% \changes{v1.2i}{2026-05-13}{Correct value of offset for large x-height}
%  \begin{macro}{\textsubscript@space}
% \changes{v1.2h}{2026-02-15}{Command added}
%    \begin{macrocode}
%</2ekernel>
%<latexrelease>\IncludeInRelease{2026/06/01}%
%<latexrelease>                 {\@textsubscript}{real text subscript}%
%<*2ekernel|latexrelease>
\protected\def\@textsubscript#1%
  {%
    \leavevmode
    \begingroup
      \check@mathfonts
      \sbox\z@{\fontsize\sf@size\sf@size#1}%
      \lower\dimexpr\textsubscript@offset\relax\box\z@
      \textsubscript@space
    \endgroup
  }
\ifdefined\directlua
  \def\textsubscript@offset{%
    \ifdim\Umathsubshiftdown\textstyle<\dimexpr\ht\z@-\Umathsubtopmax\textstyle\relax
      \ht\z@-\Umathsubtopmax\textstyle
    \else
      \Umathsubshiftdown\textstyle
    \fi
  }
\else
  \def\textsubscript@offset{%
    \ifdim\fontdimen16\textfont\tw@<\dimexpr\ht\z@-0.8\fontdimen5\textfont\tw@\relax
      \ht\z@-0.8\fontdimen5\textfont\tw@
    \else
      \fontdimen16\textfont\tw@
    \fi
  }
\fi
\def\textsubscript@space{\nobreak\hskip\scriptspace\kern\z@}
%</2ekernel|latexrelease>
%<latexrelease>\EndIncludeInRelease
%    \end{macrocode}
%
%    \begin{macrocode}
%<latexrelease>\IncludeInRelease{2020/10/01}%
%<latexrelease>                 {\@textsubscript}{real text subscript}%
%<latexrelease>\def\@textsubscript#1{%
%<latexrelease>  {\m@th\ensuremath{_{\mbox{\fontsize\sf@size\sf@size#1}}}}}
%<latexrelease>\EndIncludeInRelease
%    \end{macrocode}
%    
%    \begin{macrocode}
%<latexrelease>\IncludeInRelease{2015/01/01}%
%<latexrelease>                 {\@textsubscript}{real text subscript}%
%<latexrelease>
%<latexrelease>\def\@textsubscript#1{%
%<latexrelease>  {\m@th\ensuremath{_{\mbox{\fontsize\sf@size\z@#1}}}}}
%    \end{macrocode}
%
%    \begin{macrocode}
%<latexrelease>\EndIncludeInRelease
%<latexrelease>\IncludeInRelease{0000/00/00}%
%<latexrelease>                 {\@textsubscript}{real text subscript}%
%<latexrelease>\let\@textsubscript\@undefined
%<latexrelease>\EndIncludeInRelease
%<*2ekernel>
%    \end{macrocode}
%  \end{macro}
%  \end{macro}
%  \end{macro}
%
% \begin{macro}{\footnotesep}
%    \begin{macrocode}
\newdimen\footnotesep
%    \end{macrocode}
%
% \end{macro}
%
% \begin{macro}{\footnote}
% \changes{LaTeX2.09}{1991/11/01}
%         {(RmS) Added \cs{let}\cs{protect}\cs{noexpand} in
%          \cs{footnote}, \cs{footnotemark},
%               and \cs{footnotetext}, since \cs{xdef} is used}
% \changes{LaTeX2.09}{1991/11/22}
%         {(RmS) Added \cs{let}\cs{protect}\cs{noexpand} in
%             \cs{@xfootnote}, \cs{@xfootnotemark},
%               and \cs{@xfootnotetext}}
% \changes{LaTeX2.09}{1992/11/26}
%         {(RmS) Changed all to
%             `def`protect\string{`noexpand`protect`noexpand\string}}
% \changes{v1.1b}{1994/11/26}
%         {(ASAJ) Added \cs{protected@xdef}.}
%
%    \begin{macrocode}
\def\footnote{\@ifnextchar[\@xfootnote{\stepcounter\@mpfn
     \protected@xdef\@thefnmark{\thempfn}%
     \@footnotemark\@footnotetext}}
%    \end{macrocode}
% \end{macro}
%
% \begin{macro}{\@xfootnote}
%    \begin{macrocode}
\def\@xfootnote[#1]{%
   \begingroup
     \csname c@\@mpfn\endcsname #1\relax
     \unrestored@protected@xdef\@thefnmark{\thempfn}%
   \endgroup
   \@footnotemark\@footnotetext}
%    \end{macrocode}
% \end{macro}
%
% \begin{macro}{\@footnotetext}
% \changes{LaTeX2.09}{1991/09/29}
%     {(RmS) added \cs{reset@font}}
% \changes{LaTeX2.09}{1992/11/26}
%     {(RmS) added protection for \cs{edef}}
% \changes{v1.0a}{1994/03/07}
%     {(DPC) Extra group for color}
% \changes{v1.0c}{1994/03/14}
%     {(DPC) Use \cs{color@begingroup}, add \cs{endgraf}}
% \changes{v1.0d}{1994/04/18}
%     {(DPC) Remove Color support}
% \changes{v1.0g}{1994/05/13}
%     {(DPC) Add new style color support: \cs{normalcolor}}
% \changes{v1.0g}{1994/05/13}
%     {(DPC) Use \cs{@finalstrut}}
% \changes{v1.1a}{1994/10/31}
%     {(DPC/CAR) Move color setting to output routine}
% \changes{v1.1b}{1994/11/04}
%     {(ASAJ) Added \cs{protected@edef}.}
% \changes{v1.1c}{1994/11/05}
%         {Removed \cs{normalcolor} (again)}
% \changes{v1.1t}{1997/11/19}
%         {Missing percent, again}
% \changes{v1.2e}{2021/02/10}{Explicitly run \cs{par} at the end of footnote text
%      in preparation for paragraph hooks}
% \changes{v1.2g}{2021/10/14}
%         {Explicitly set \cs{@currentcounter} (gh/687)}
%    \begin{macrocode}
%</2ekernel>
%<*2ekernel|latexrelease>
%<latexrelease>\IncludeInRelease{2021/11/15}%
%<latexrelease>                 {\@footnotetext}{footnotetext tagging}%
\long\def\@footnotetext#1{\insert\footins{%
    \reset@font\footnotesize
    \interlinepenalty\interfootnotelinepenalty
    \splittopskip\footnotesep
    \splitmaxdepth \dp\strutbox \floatingpenalty \@MM
    \hsize\columnwidth \@parboxrestore
    \def\@currentcounter{footnote}%
    \protected@edef\@currentlabel{%
       \csname p@footnote\endcsname\@thefnmark
    }%
    \color@begingroup
      \@makefntext{%
        \rule\z@\footnotesep\ignorespaces#1\@finalstrut\strutbox}%
    \par
    \color@endgroup}}%
%</2ekernel|latexrelease>
%<latexrelease>\EndIncludeInRelease
%    \end{macrocode}
%
%    \begin{macrocode}
%<latexrelease>\IncludeInRelease{2021/06/01}%
%<latexrelease>                 {\@footnotetext}{footnotetext tagging}%
%<latexrelease>\long\def\@footnotetext#1{\insert\footins{%
%<latexrelease>    \reset@font\footnotesize
%<latexrelease>    \interlinepenalty\interfootnotelinepenalty
%<latexrelease>    \splittopskip\footnotesep
%<latexrelease>    \splitmaxdepth \dp\strutbox \floatingpenalty \@MM
%<latexrelease>    \hsize\columnwidth \@parboxrestore
%<latexrelease>    \protected@edef\@currentlabel{%
%<latexrelease>       \csname p@footnote\endcsname\@thefnmark
%<latexrelease>    }%
%<latexrelease>    \color@begingroup
%<latexrelease>      \@makefntext{%
%<latexrelease>        \rule\z@\footnotesep\ignorespaces#1\@finalstrut\strutbox}%
%<latexrelease>    \par
%<latexrelease>    \color@endgroup}}%
%<latexrelease>\EndIncludeInRelease
%    \end{macrocode}
%    
%    \begin{macrocode}
%<latexrelease>\IncludeInRelease{0000/00/00}%
%<latexrelease>                 {\@footnotetext}{footnotetext tagging}%
%<latexrelease>
%<latexrelease>\long\def\@footnotetext#1{\insert\footins{%
%<latexrelease>    \reset@font\footnotesize
%<latexrelease>    \interlinepenalty\interfootnotelinepenalty
%<latexrelease>    \splittopskip\footnotesep
%<latexrelease>    \splitmaxdepth \dp\strutbox \floatingpenalty \@MM
%<latexrelease>    \hsize\columnwidth \@parboxrestore
%<latexrelease>    \protected@edef\@currentlabel{%
%<latexrelease>       \csname p@footnote\endcsname\@thefnmark
%<latexrelease>    }%
%<latexrelease>    \color@begingroup
%<latexrelease>      \@makefntext{%
%<latexrelease>        \rule\z@\footnotesep\ignorespaces#1\@finalstrut\strutbox}%
%<latexrelease>    \color@endgroup}}%
%<latexrelease>
%<latexrelease>\EndIncludeInRelease
%<*2ekernel>
%    \end{macrocode}
% \end{macro}
%
% \begin{macro}{\footnotemark}
% \changes{v1.1b}{1994/11/04}{Added \cs{protected@xdef} to
%    \cs{footnotemark}.}
%    \begin{macrocode}
\def\footnotemark{%
   \@ifnextchar[\@xfootnotemark
     {\stepcounter{footnote}%
      \protected@xdef\@thefnmark{\thefootnote}%
      \@footnotemark}}
%    \end{macrocode}
% \end{macro}
%
% \begin{macro}{\@xfootnotemark}
%    \begin{macrocode}
\def\@xfootnotemark[#1]{%
   \begingroup
      \c@footnote #1\relax
      \unrestored@protected@xdef\@thefnmark{\thefootnote}%
   \endgroup
   \@footnotemark}
%    \end{macrocode}
% \end{macro}
%
% \begin{macro}{\@footnotemark}
% \changes{v1.1h}{1995/05/12}
%         {Add \cs{nobreak} to allow hyphenation. latex/1605}
%    \begin{macrocode}
\def\@footnotemark{%
  \leavevmode
  \ifhmode\edef\@x@sf{\the\spacefactor}\nobreak\fi
  \@makefnmark
  \ifhmode\spacefactor\@x@sf\fi
  \relax}
%    \end{macrocode}
% \end{macro}
%
% \begin{macro}{\footnotetext}
%    \begin{macrocode}
\def\footnotetext{%
     \@ifnextchar [\@xfootnotenext
       {\protected@xdef\@thefnmark{\thempfn}%
    \@footnotetext}}
%    \end{macrocode}
% \end{macro}
%
% \begin{macro}{\@xfootnotenext}
%    \begin{macrocode}
\def\@xfootnotenext[#1]{%
  \begingroup
     \csname c@\@mpfn\endcsname #1\relax
     \unrestored@protected@xdef\@thefnmark{\thempfn}%
  \endgroup
  \@footnotetext}
%    \end{macrocode}
% \end{macro}
%
% \begin{macro}{\thempfn}
% \begin{macro}{\@mpfn}
%    \begin{macrocode}
\def\@mpfn{footnote}
\def\thempfn{\thefootnote}
%    \end{macrocode}
%
% \end{macro}
% \end{macro}
%
%
%
%
%  \begin{macro}{\footref}
%    This command generates a footnote mark. The value is produced by
%    referencing a \cs{label} placed into a \cs{footnote} elsewhere
%    (can be one in the main galley or in a minipage).
% \changes{v1.2f}{2021/02/16}{\cs{footref} added}
%    \begin{macrocode}
%</2ekernel>
%<*2ekernel|latexrelease>
%<latexrelease>\IncludeInRelease{2021/06/01}%
%<latexrelease>                 {\footref}{Add footref}%
\def\footref#1{%
  \begingroup
    \unrestored@protected@xdef\@thefnmark{\ref{#1}}%
  \endgroup
  \@footnotemark
}
%</2ekernel|latexrelease>
%<latexrelease>\EndIncludeInRelease
%    \end{macrocode}
%    We don't remove it when rolling back so that packages offered it
%    in the past do not need to alter their behavior in a rollback
%    situation.
%    \begin{macrocode}
%<latexrelease>\IncludeInRelease{0000/00/00}%
%<latexrelease>                 {\footref}{Add footref}%
%<latexrelease>
%<latexrelease>  % \let\footref\@undefined
%<latexrelease>
%<latexrelease>\EndIncludeInRelease
%<*2ekernel>
%    \end{macrocode}
%  \end{macro}
%
%
%    \begin{macrocode}
%</2ekernel>
%    \end{macrocode}
%
% \Finale
%
